\documentclass[11pt]{article}
\usepackage{amsmath,amssymb,amsthm}
\usepackage[margin=1.1in]{geometry}
\usepackage{hyperref}

\newtheorem{theorem}{Theorem}[section]
\newtheorem{lemma}[theorem]{Lemma}
\newtheorem{proposition}[theorem]{Proposition}
\newtheorem{corollary}[theorem]{Corollary}
\theoremstyle{definition}
\newtheorem{definition}[theorem]{Definition}
\theoremstyle{remark}
\newtheorem{remark}[theorem]{Remark}
\newtheorem{openq}[theorem]{Open problem}

\newcommand{\pmS}{\operatorname{pm}}
\newcommand{\amp}{\operatorname{amp}}

\title{A machine-checked programme for the Krenn--Gu conjecture}
\author{}
\date{}

\begin{document}
\maketitle

\begin{abstract}
Every lemma stated here is verified by the Lean kernel, on \texttt{propext},
\texttt{Classical.choice} and \texttt{Quot.sound} and nothing further. The
conjecture itself is \emph{not} proved. What is offered is an architecture, the
lemmas so far established within it, and --- stated as precisely as the proved
material allows --- the obstructions that remain. Two claims made during the work
are recorded as withdrawn, with the reasons, since a programme that hides its wrong
turns is worth less than one that keeps them.
\end{abstract}

\section{The statement}

Let $V$ be a finite vertex set of even size and $D$ a number of colours. A
\emph{weight system} assigns to each unordered pair of coloured vertices a complex
number. Its \emph{amplitude} at a colouring $c\colon V\to D$ is
\[
  \amp_W(c)\;=\;\sum_{\sigma}\ \prod_{e\in\sigma} W\bigl(e\text{ coloured by }c\bigr),
\]
the sum over perfect matchings of $V$. The system \emph{solves} the equations if
$\amp_W(c)=1$ for every constant $c$ and $0$ for every other $c$.

\begin{remark}
The amplitude is the \emph{hafnian} of the coloured weight matrix. Its signed
counterpart, the Pfaffian, squares to a determinant, evaluates in polynomial time,
and satisfies contraction identities relating a matrix to smaller ones through its
inverse. The hafnian has none of these. This is not an analogy: it is a structural
reason to expect arguments of contraction type to fail here, and they do
(\S\ref{sec:withdrawn}).
\end{remark}

\section{What is proved}

\subsection{The problem is correctly posed}

\begin{theorem}[Sharpness at four vertices]
Four vertices carry a solution in three colours. Explicitly: their three perfect
matchings are the three nonzero classes of the Klein four-group structure on $V$;
give each matching a colour and weight each edge $1$ exactly when both endpoints
wear that edge's colour.
\end{theorem}

\begin{theorem}[The bound of two is attained]
At every even vertex count there is a solution in two colours: colour the edges of a
Hamiltonian cycle by which of its two alternating matchings they belong to, and
leave the chords unweighted.
\end{theorem}

\noindent The second rests on the fact that an even cycle has exactly two perfect
matchings, proved here by an alternation argument over $\mathbb{Z}/m$.

\subsection{Gauge freedom}

\begin{theorem}[Normalisation is free]
Multiplying each weight by factors attached to its endpoints' colours rescales the
amplitude by a single number, uniformly across matchings, since every matching
covers every vertex once. Hence a system whose monochromatic amplitudes are merely
nonzero rescales to one where they equal $1$, with all vanishing conditions
preserved.
\end{theorem}

\subsection{The vertex equations}

Fix $u\in V$. Expanding the amplitude along $u$,
\[
  \amp_W(c)\;=\;\sum_{v\neq u} W\bigl(uv\text{ coloured by }c\bigr)\cdot
  \pmS_W\bigl(c,\;V\setminus\{u,v\}\bigr),
\]
where $\pmS_W(c,S)$ is the matching sum over the subset $S$. The unknowns are
indexed by $u$'s \emph{partners}: the system's size is $u$'s degree, and nothing in
it grows with $|V|$.

\begin{theorem}[Constant background]
For each $k$ and each colour $j$ that $u$ may wear,
$\sum_{v\neq u} W(u{:}j,\,v{:}k)\,\pmS_W(\mathbf{k},V\setminus\{u,v\})
 = [\,j=k\,]$.
\end{theorem}

\begin{theorem}[Non-constant background]
If $c$ is already non-constant away from $u$, the right-hand side is $0$ for
\emph{all three} $j$. This requires no normalisation, only that non-monochromatic
colourings vanish.
\end{theorem}

\begin{theorem}[Independence of the monochromatic row]
At every vertex and colour, the weights keeping that vertex monochromatic are not a
linear combination of those breaking it. No vector is annihilated by two rows and
sent to $1$ by their combination.
\end{theorem}

\begin{proposition}[Low degree]
A vertex with one partner carries only matching colours on that edge. A vertex with
two forces an explicit $2\times2$ determinant of colour-breaking weights to vanish.
\end{proposition}

\subsection{A case of the conjecture, from cubic graph theory}

Call a weight system \emph{colour-per-matching} when each edge carries weight at a
single colour. The edges of each colour then form a perfect matching, each matching
wears one colour, and this is precisely the construction giving the four-vertex
three-colour solution above. It is also the branch that the structural split of
\S\ref{sec:obstruct} isolates: solutions whose edges are single-coloured are of this
kind.

\begin{theorem}[The colour-per-matching case]\label{thm:cpm}
For three colours no colour-per-matching solution exists on more than four
vertices.
\end{theorem}

\begin{proof}[Argument]
Such a solution needs three edge-disjoint perfect matchings whose union $U$ has
\emph{exactly} those three: a further matching of $U$ would mix colours, and a
colouring constant on each of its edges but not overall would have nonzero
amplitude.

$U$ is cubic and carries a proper three-edge-colouring, hence is bridgeless, since a
cubic graph with a bridge admits no such colouring.

The union of two of the classes is a disjoint union of even cycles covering every
vertex. Were there two or more cycles, exchanging along one of them alone produces a
fourth perfect matching --- inside the exchanged set it follows the second class
where the first disagreed, outside it follows the first where the second disagreed.
So every pair of classes forms a single Hamiltonian cycle, and the third class is a
perfect matching of chords of that cycle.

Now let $\{u,v\}$ be a chord whose removal splits the cycle into two paths with
evenly many vertices each. That chord together with the cycle edges along both paths
is a perfect matching: it contains a third-class edge, so is neither of the first
two, and contains cycle edges, so is not the third. A fourth matching again.

On four vertices the chords are the diagonals of a four-cycle, and deleting a
diagonal's endpoints leaves two \emph{single} vertices --- odd paths, no fourth
matching. That is the whole of the exception. From six vertices upward a chord of
even span exists.
\end{proof}

\begin{remark}
Exhaustively over all three-edge-coloured cubic subgraphs of the complete graph:
one such subgraph on four vertices and it has exactly three matchings; $80$ on six
vertices and none does; $32\,970$ on eight vertices and none does. The six-vertex
witness is concrete: take the six-cycle as the union of the first two classes and
the three diameters as the third; removing a diameter leaves two paths of two
vertices, so that diameter with four cycle edges is a fourth matching. This is the
same configuration an independent exact test locates as the colouring on which the
construction's amplitude fails.
\end{remark}

\noindent The exchange step is machine-checked; the chord-parity step is not yet.

\subsection{The vertex equations as a matrix identity}

Write $A^{ij}$ for the symmetric matrix of weights at colour pair $(i,j)$, so that
the amplitude at a colouring $c$ is the hafnian of the matrix whose $(u,v)$ entry is
$A^{c_uc_v}_{uv}$. Write
\[
  H^{(k)}_{uv} \;=\; \operatorname{haf}\bigl(A^{kk}\restriction V\setminus\{u,v\}\bigr)
\]
for the \emph{hafnian minors} of the monochromatic matrix. Expanding along a vertex
turns the equations into row sums of Hadamard products:
\[
  \bigl(A^{kk}\circ H^{(k)}\bigr)\mathbf{1} \;=\; \mathbf{1},
  \qquad
  \bigl(A^{mk}\circ H^{(k)}\bigr)\mathbf{1} \;=\; \mathbf{0}
  \quad (m\neq k).
\]

The first is empty: it is the hafnian expansion of $\operatorname{haf}A^{kk}=1$ along
each row. The content is the second --- \emph{each off-diagonal weight matrix is
row-orthogonal to the hafnian minors of the diagonal one}.

For determinants the analogous object is the adjugate, and $A\operatorname{adj}A =
(\det A)I$ pins it down completely. There is no hafnian counterpart. That absence is
the hafnian/Pfaffian gap in its sharpest form: the minors $H^{(k)}$ are as hard to
compute as the hafnian itself and satisfy no identity relating them back to
$A^{kk}$. The equations therefore constrain a quantity that the theory offers no
handle on, and this is why a determinantal or contraction-style argument has nothing
to grip.

\subsection{Compatible matchings, and one recurring invariant}

Call a matching \emph{compatible} with a colouring when each of its edges carries
nonzero weight at the colours that colouring gives the edge's endpoints. Only
compatible matchings contribute.

\begin{theorem}[No unique compatible matching]\label{thm:uniq}
No non-constant colouring has exactly one compatible matching.
\end{theorem}

\begin{proof}
Its product would be the whole amplitude, which must vanish; a product of nonzero
factors does not.
\end{proof}

\noindent Theorem~\ref{thm:uniq} mentions no complex number: it constrains the
\emph{supports} alone, and refuting a candidate support system needs only one
non-constant colouring admitting exactly one matching. It is Theorem~\ref{thm:cpm}
stripped of its graph-theoretic clothing --- there the mixed colouring admits
precisely one matching.

\begin{theorem}[Two compatible matchings]
If a non-constant colouring admits exactly two compatible matchings, their products
are negatives of one another.
\end{theorem}

\noindent This is the refinement that reaches where Theorem~\ref{thm:uniq} cannot.
That theorem yields a \emph{contradiction}, but only when a colouring singles out
one matching, which never happens at full support. This yields an \emph{equation},
and equations survive into the dense regime.

\begin{lemma}[Localisation]
If two matchings agree outside a set closed under both, and the product over the
complement does not vanish, a relation between their whole products is the same
relation between their products over the set alone.
\end{lemma}

\noindent Three apparently separate results are now one. A colour-slice of size two
forces a $2\times2$ determinant to vanish; a local family cancels when its products
sum to zero; two compatible matchings differing by one alternating four-cycle have
opposite products. Localised, each says
\[
  \pmS(A) = 0 \quad\text{for a four-vertex set } A.
\]
The recurring invariant is the matching sum over a small set, and the equations
reach it three different ways.

\begin{theorem}[Factorisation dichotomy]
If a colouring kills every matching crossing the boundary of a set, then the
matching sum over that set or over its complement vanishes.
\end{theorem}

\noindent What remains factors as the two sides' matching sums, the amplitude
vanishes, and $\mathbb{C}$ has no zero divisors. \emph{Which} side vanishes is not
determined --- which is exactly why this is a dichotomy and not a descent.

\subsection{Where the combinatorics stops}

Measured at six vertices and three colours over random supports, each edge keeping
each colour pair with probability $p$, Theorem~\ref{thm:uniq} alone refutes:
\[
\begin{array}{llllllll}
p: & 0.15 & 0.25 & 0.35 & 0.5 & 0.7 & 0.9 & 1.0\\
\text{refuted}: & 92\% & 99\% & 100\% & 100\% & 94\% & 3\% & 0\%
\end{array}
\]
The method covers everything except near-full support and has exactly zero purchase
at full support, where every matching is compatible with everything. The transition
is sharp.

This is the third independent route to the same place. The support abstraction is
vacuous at full support; the exact algebra needs some hundred and twenty live
variables against a tractability ceiling near fifty; and now the combinatorics dies
there too. The three are not separate obstacles. What makes a support dense enough
to defeat the combinatorics is what makes its algebra large.

(These are \emph{random} supports. A counterexample's support would be structured,
and structure sits anywhere. The figures bound the method's reach, not the
probability of catching a counterexample.)

\subsection{Where the missing identity does exist}

The obstruction of the previous subsection is the absence of a hafnian adjugate. It
is worth asking where that absence is repaired.

A graph admits a \emph{Pfaffian orientation} when signs can be attached to its edges
so that every perfect matching acquires the same sign. Kasteleyn's theorem supplies
one for every planar graph, and Little's for every graph without a $K_{3,3}$ minor
in the bipartite case. Where such an orientation exists the hafnian \emph{is} a
Pfaffian up to sign, $\operatorname{Pf}(A)^2=\det(A)$ holds, and the adjugate
identity returns with all of determinantal theory behind it.

On that side of the divide the defining equations become determinantal. The
monochromatic conditions read $\det = 1$ and every mixed colouring reads $\det = 0$;
and a colouring differing from a constant one at a single vertex modifies exactly one
row and column, so the vanishing is a rank statement about a structured perturbation
of a matrix of determinant one. That is a problem with tools.

This suggests the axis along which the remaining case should be split is not density
but \emph{Pfaffian orientability}. The two are related --- complete graphs on six or
more vertices are as far from Pfaffian as graphs get, and full support is exactly the
complete graph --- which would explain, rather than merely restate, why every method
tried here fails in the same place.

\begin{proposition}[Verified]
A uniform-sign orientation exists for $K_4$ and for no larger complete graph:
\[
  K_4:\ \text{yes}\qquad K_6:\ \text{no}\qquad K_8:\ \text{no}.
\]
\end{proposition}

\begin{proof}[Method]
Such an orientation exists exactly when the system over $\mathbb{F}_2$ requiring each
matching's edge-parity to differ from its permutation sign by a constant is solvable.
Gaussian elimination settles each case exactly.
\end{proof}

\noindent Since the vertex count is even and $K_4$ is the largest planar complete
graph, $K_4$ is the \emph{unique} even complete graph that is Pfaffian. It is also
the exceptional case of the conjecture. \textbf{The vertex count at which a solution
exists is the vertex count at which the missing identity exists.}

This is the second independent account of that exceptionality. The first, from cubic
graph theory, made $K_4$ the unique cubic graph all of whose chords have odd span
(Theorem~\ref{thm:cpm}). Different subject, different argument, same object singled
out.

None of this proves the conjecture: a graph failing to be Pfaffian does not thereby
fail to carry a solution. What it does is motivate the axis. Full support \emph{is}
the complete graph, and the complete graph on six or more vertices is exactly where
the Pfaffian property is lost --- so if the real divide is orientability, then
``dense'' was a proxy for ``non-Pfaffian'' throughout, and the three failures
recorded in \S\ref{sec:obstruct} are one failure with one cause.

\begin{remark}[The converse fails]
Being Pfaffian is not sufficient for solvability. On six vertices both the cycle and
the prism are planar, hence Pfaffian, and a numerical search finds no three-colour
solution on either --- as it finds none on the non-Pfaffian $K_{3,3}$. So the
property cannot be the criterion for carrying a solution; at most it is a necessary
ingredient, and the coincidence at $K_4$ is evidence about the \emph{obstruction},
not about sufficiency. (Three restarts each; weak evidence, recorded as such.)
\end{remark}

Formalising Kasteleyn's theorem is a large undertaking and was not attempted.

\section{The cap framework, machine-checked}

An external structural review supplied a sharper reduction than anything above: the
conjecture follows from one \emph{integrable-cap} theorem. Its own ledger recorded
that none of its results had been checked by Lean. The following now have been, each
depending on \texttt{propext}, \texttt{Classical.choice}, \texttt{Quot.sound} and
nothing more.

\begin{center}
\begin{tabular}{ll}
\hline
Colour restriction $d\to3$ & checked\\
Cap identity: the pair split & checked\\
Constant-diagonal cap of the target & checked\\
Integrability: no matching uses two defect edges & checked\\
Integrability: a one-factor perturbation is exactly linear & checked\\
Absorption: defect of matching number one, any shape & checked\\
The $n\to n-2$ reduction & arithmetic remaining\\
\emph{Cap existence} & \emph{open --- the difficulty}\\
\hline
\end{tabular}
\end{center}

Two observations are worth recording.

First, the framework retro-explains the failure of \S\ref{sec:obstruct}. Contracting
two sites leaves a flat object \emph{plus a first-order defect}, and the defect is not
itself flat. The fragment-descent route abandoned earlier was demanding that the
defect \emph{vanish}; the framework's insight is that it need only be
\emph{absorbed}, which is possible exactly when its edges have matching number at
most one. Asking for too much is why that route saw only the low-connectivity regime.

The absorption holds for every shape the condition permits, not only the star. What
made the general case go through was writing the first-order part as a sum over the
whole defect family rather than as a case split: the no-defect and one-defect cases
then become the same expression, and the shape of the support never enters the
proof. The classification into star or triangle, which the framework offers as an
aid to intuition, turns out to be unnecessary --- pairwise intersection is what the
argument uses.

Second, almost nothing had to be built. The deletion bijection --- matchings sending
one vertex to another are the reinstatements of matchings of the rest, by right
multiplication with a transposition --- was proved for an unrelated purpose, as were
amplitude naturality, locality of the subset matching sum, and the factorisation of an
edge product across a closed set. That the cap framework's combinatorial content is
expressible almost entirely in machinery built independently is evidence both that
these abstractions are the right ones and that the framework is natural rather than
contrived.

\subsection{Where cap existence reaches}

The framework's open assertion --- that every hypothetical solution admits a cap
whose defect has matching number at most one --- is \emph{recognisable} by cheap
linear algebra. For a vertex pair and a candidate defect shape, the conditions
``the defect vanishes off that shape'' are linear in the cap; solve, take the null
space, and ask for a vector with nonzero diagonal pairing nontrivially with the
contracted edge. Running that recognizer:

\begin{center}
\begin{tabular}{lr}
\hline
six-vertex two-colour solution (proved, sparse) & $648$ caps\\
random dense, six vertices, three colours & $0$ of $10$\\
random dense, eight vertices, three colours & $0$ of $10$\\
\hline
\end{tabular}
\end{center}

Both permitted shapes were searched, star and triangle, which on a simple graph are
the only ones with matching number at most one. Caps are abundant for the sparse
solution --- its found defects had support of size one, so almost any cap works
without anything being arranged --- and absent for dense systems.

Dense systems are not solutions, so this refutes nothing. What it rules out is the
cheapest route to proving cap existence:

\begin{quote}
Cap existence, if true, must be forced by the defining equations. It cannot follow
from genericity or dimension counting, since a generic system has no cap.
\end{quote}

The counting is stark. For three colours the constant-diagonal caps form a
seven-dimensional space; confining the defect to a star at eight vertices kills
about ten vertex pairs at nine scalar conditions each. Ninety equations against
seven unknowns. Only a strong degeneracy satisfies that, and the sole available
source of degeneracy is the solution's own equations.

This is the fourth method to fail at dense support, after the support abstraction,
the exact algebra, and the exclusion of a unique compatible matching. Either the
boundary is the real object --- dense support \emph{is} the complete graph, where the
Pfaffian property is lost --- or every method so far has been one idea in disguise.

\section{Contraction and the incidence matrix}\label{sec:contract}

Contracting a pair of sites replaces the amplitude by a flat term together with a
correction. The corrections leaving the amplitude flat away from a permitted defect
shape are the kernel of a matrix assembled from the weights.

\begin{lemma}[Support restriction; sufficient, not necessary]
An entry of that matrix vanishes whenever the weight block of a meeting edge is
identically zero. No hypothesis on the origin of the weights is needed. The converse
fails: an entry may vanish by cancellation inside the cap matrix with every block
nonzero, so failure of this criterion does not establish that no cap exists.
\end{lemma}

Call a residual pair \emph{live} when one of its vertices adjoins each contracted
site. Only live pairs can contribute a constraint.

\begin{proposition}
A family of pairs lies inside one star or one triangle if and only if it is pairwise
intersecting. The restriction to those two shapes is therefore forced, not chosen.
\end{proposition}

\begin{lemma}[Star collapse]
If a contracted site meets the residual set in at most one place, every live pair
contains that place, the family is a star, and the system is empty.
\end{lemma}

Whether a \emph{usable} cap exists is not settled by the shape alone. Two linear
functionals on the kernel must both avoid zero: the constant diagonal and the
contraction scalar.

\begin{lemma}[Two functionals, one witness]
If each of two linear functionals is nonzero somewhere on a subspace, some single
vector of that subspace makes both nonzero, and it is one of $u$, $v$, $u+v$.
\end{lemma}

Consequently a search that inspects basis vectors individually answers a strictly
weaker question and under-reports cap existence.

\begin{remark}[Caps at eight vertices are non-generic]
On the cube, contracting a pair with a star defect admits a cap with constant nonzero
diagonal $\kappa=-3/4$, contraction scalar $-91/16$, and all ninety flatness
constraints satisfied exactly over the rationals. Such caps are not generic: with
wide-range weights none appear on the cube or on Wagner's graph, and with small integer
weights they appear exactly when some edge block is singular.

Each incidence row's coefficient array is a symmetrised product $u(k,x)\otimes v(l,y)$
of site vectors $u(k,x)_a=W_{(p,a),(k,x)}$ and $v(l,y)_b=W_{(q,b),(l,y)}$. On a cubic
support only neighbours of the contracted sites contribute, and the surviving rows span
a tensor product of two site spans, filling the cap space unless a factor drops rank.
Cap existence at these sizes is therefore a determinantal condition on the weights.
\end{remark}

\begin{openq}
Does the GHZ variety meet the locus on which an edge block is singular? If it does, the
contraction route reaches eight vertices; if not, the route cannot reach the first open
case, for a reason resting on the defining equations rather than on any recognizer.
Sampling cannot decide this, the relevant points being exactly the non-generic ones.
\end{openq}

\section{Second-order expansion}\label{sec:second}

A defect on two edges meets any matching in at most two of them, so the perturbed
matching sum is an exact quadratic: a flat term, a linear term, and a cross term,
arising when one matching uses both defect edges.

\begin{theorem}[Second-order absorption]
The rewriting that folds a first-order defect into perturbed weights still applies,
leaving a residue carrying the quadratic coefficient. The reduction returns an object
of the same kind exactly when that coefficient vanishes.
\end{theorem}

\begin{theorem}[The coefficient factors]
The quadratic coefficient is the defect's values at its two edges times a double
cofactor depending only on the weights. For disjoint defect edges that cofactor is the
matching sum over the vertices they leave behind.
\end{theorem}

This development stands on its own; it is not forced by any proved failure of the
first-order route. The four-vertex matching sum it meets recurs across several
arguments -- the factorisation dichotomy, the criterion for a local family to cancel,
the determinant at a colour slice of degree two, the exceptional standing of four
vertices -- but it is a candidate crux rather than an object shown to control every
route. What is known about it is a dichotomy: one side of a partition has vanishing
matching sum, with which side left open, and that openness is a property of the
statement rather than a defect in its proof.

\section{What the equations force on a block}\label{sec:blockrank}

Fix a colouring that already disagrees somewhere away from two chosen sites. Then the
colouring stays non-constant however those two sites are painted, so a solution's
matching sum vanishes identically in their colours. Splitting matchings by how the two
sites are matched gives an identity, valid for arbitrary weights:

\begin{theorem}[Block decomposition]
With the colours at two sites free and every other colour fixed, the matching sum
equals the block those sites carry, scaled by the matching sum on the remaining sites,
plus one rank-one term for each ordered pair of distinct partners the two sites reach.
\end{theorem}

Grouping by the outer index shows the number of rank-one terms is the number of
partners the \emph{first} site reaches, not the number of pairs. A site of degree three
therefore contributes two.

\begin{theorem}[Degeneration]
If a square matrix has every entry a sum of fewer row-column products than it has rows,
some nonzero combination of its rows vanishes.
\end{theorem}

\begin{theorem}[Dichotomy]
Let a solution have a site $p$ reaching fewer partners than there are colours, and let
$q$ be any other site. Then either the block $W_{pq}$ has linearly dependent rows, or
every colouring disagreeing away from $p,q$ annihilates the matching sum on the
remaining sites --- the vanishing half of the defining conditions two sites down.
\end{theorem}

The second alternative needs a nonvanishing monochromatic sum to become a solution on
two fewer sites, and the equations supply one:

\begin{theorem}[Live partners]
For every site and every colour, some partner has both a nonzero block entry and a
nonzero monochromatic matching sum on the remaining sites.
\end{theorem}

This is read off the vertex equation, which at equal colours sums to one; a sum equal to
one cannot have every summand vanish.

\section{A route that must fail}\label{sec:mustfail}

The descent would close immediately if one partner were live for every colour at once.
On weights built by colouring the matchings that never happens: such an edge wears a
single colour and carries weight only there, so it is live for one colour only. The
four-vertex solution is of this kind, and there each colour has its own live partner,
the three being distinct.

This does not show a simultaneous witness never exists --- a hypothetical solution need
not come from a colouring --- but it does show one cannot be assumed. And the failure is
compulsory rather than unfortunate: four vertices are solvable and must not descend, so
every correct descent argument has to fail there. A route through a simultaneous witness
fails there for the wrong reason, and would prove too much.

The numerology is worth noting. At four vertices a site has three partners against three
colours and the assignment of colours to live partners is a bijection --- the same
coincidence that makes four vertices exceptional. A cubic support on eight vertices
reproduces it: three neighbours, three colours. Any descent must break that coincidence
rather than count around it.

\section{What blocks a cap}\label{sec:capblock}

The reduction needs a cap whose correction is absorbable. The most natural construction
is an outer product, and it fails for reasons the equations themselves supply.

\begin{theorem}[No annihilator at a site]
No nonzero vector annihilates every weight block at a site. The vertex equation exhibits
each basis vector as a combination of the blocks' columns in the corresponding colour, so
those columns span, and pairing a candidate annihilator against the equation returns that
vector's coordinate.
\end{theorem}

\begin{corollary}
An outer-product cap is never flat: its first factor would have to annihilate every block
at the site.
\end{corollary}

A cap that is flat only off a permitted shape fares no better in one direction:

\begin{theorem}[Self-defeat]
An outer-product cap contracts to the pairing of its second factor against the first
applied to the block's columns. A first factor annihilating the block therefore leaves no
contraction scalar. Where the partners are few, the equations force annihilating the
partners to imply annihilating the block, so the two requirements cannot both hold.
\end{theorem}

So a usable cap, if one exists, is a genuinely higher-rank matrix rather than an outer
product. Whether it exists is unresolved and asserted nowhere.

\section{How far the degeneration goes}\label{sec:howfar}

The dichotomy forces a block at a low-degree site to have dependent rows. It forces more
than that.

\begin{theorem}[Rank two is excluded]
A block of rank two at such a site would have a line for its annihilator, so the
annihilators produced at different colourings would all point one way; that single
direction would then annihilate every block at the site, hence vanish. So the block has
rank at most one.
\end{theorem}

Four vertices escapes precisely here. Its blocks have rank one, the annihilator is a
plane rather than a line, the direction is not pinned, and different colourings may use
different ones. Since four vertices is solvable and must not descend, any correct argument
has to fail there --- and this one fails for exactly the right reason.

Under rank one the blocks factor, and the vertex equation says each basis vector is a
combination of the blocks' first factors. On a support where a site has three partners
those three factors are then a basis, which is the same numerical coincidence that makes
four vertices exceptional. Whether it can be sustained above four vertices is open.

\section{The scope of the degeneration chain}\label{sec:scope}

The chain was first stated at a site with fewer partners than there are colours, but the
partner count is not what it depends on. A correction indexed by partners reindexes along
any spanning family for the vectors those partners contribute, so the number of terms can
be taken to be the dimension those vectors span. Degree is irrelevant; dimension is what
matters.

\begin{remark}[The actual hypothesis]
The degeneration runs exactly when the columns a colouring uses at a site --- one per
partner, at the colour that colouring assigns --- span fewer dimensions than there are
colours. Nothing in the equations forces this. With three colours and four or more
partners, columns in general position span all three, and then the identity reads that a
three-by-three block equals a correction of rank at most three, which says nothing.
\end{remark}

So the chain is conditional on a rank drop among the used columns, and the condition is
not vacuous but neither is it supplied. This is the honest position: the degeneration
argument is available on the locus where the used columns degenerate, and the equations do
not put a solution on that locus.

The self-limitation is now visible in its proper form. The vertex equation makes the
columns, taken over all colours, span everything. A colouring lowers the span only by
using columns that happen to be dependent --- and under rank-one factorisation that means
using colours where second factors vanish. The chain's own conclusion is that a witnessing
colouring must do exactly that, so the argument establishes the condition it needs rather
than deriving a contradiction from its failure.

What would close it is a proof that some colouring is forced onto a degenerate set of
columns, or a construction showing none need be. Neither is here.

\section{The site relation}\label{sec:site}

The contraction framework needed a cap, and a cap needed the blocks to degenerate, and
the argument meant to supply that degeneracy could only conclude the condition it
required. The way out is to obtain degeneracy from somewhere else. Expanding the
amplitude at a site does exactly that, and asks nothing of the framework.

\begin{theorem}[The site relation]
Expanding the amplitude at a site leaves one term for each neighbour it reaches. At a
colouring already non-constant away from that site, a solution makes those terms cancel
--- and they cancel for every colour the site itself wears, since no coefficient sees it.
So the columns the neighbours contribute are linearly dependent, with the complementary
matching sums as coefficients.
\end{theorem}

With three colours the statement is sharp at degree three: three columns in three
dimensions, dependent exactly when their determinant vanishes.

\begin{corollary}
If those columns are independent, every complementary matching sum vanishes. Colouring
everything one way but a chosen neighbour turns this into the vanishing of a
\emph{monochromatic} matching sum on the complement of the site and that neighbour.
\end{corollary}

The vertex equation says some neighbour keeps such a sum alive at each colour, and a site
reaches no neighbours beyond its own. So independence cannot hold everywhere, and reading
the implication backwards:

\begin{theorem}[Liveness forces degeneracy]
Wherever a colouring non-constant away from the site leaves a complementary matching sum
alive, the columns that colouring selects are dependent.
\end{theorem}

The determinantal conditions are therefore indexed by live colourings, not merely by the
three the vertex equation supplies directly.

\section{What the degeneracy gives, and what it does not}\label{sec:sitelimit}

Two pieces of linear algebra convert these conditions into geometry. A dependent triple
whose last two members are independent puts the first in their span, since its coefficient
cannot vanish. And no plane can hold every column at a site: two vectors in three colours
always admit a common annihilator, and an annihilator of every column at a site is zero.

\begin{corollary}
If the planes spanned by the pairs of columns at the other neighbours coincide across
colours, every column at the site lies in that one plane, which is impossible.
\end{corollary}

That closes one case. In the other, the planes are distinct and each column of the live
neighbour's block is confined to the line where two of them meet --- a more rigid
situation, not a weaker one, but not yet a contradiction.

What decides it is how many colourings are live. Dependency is forced at each of them, so
a solution evades the argument only by having complementary sums vanish at nearly every
colouring. That is a statement about matching sums on two fewer sites, and it is where the
cubic case now stands: not on a question of geometry, but on a count.

\section{The spread case}\label{sec:spread}

The descent branch needs one neighbour whose removal leaves all three monochromatic
matching sums alive. The equations supply one live neighbour per colour; whether a single
neighbour serves all three is the question. Where none does --- three colours served by
three different neighbours --- the configuration is rigid rather than free.

\begin{theorem}[Pinning]
If at some colour exactly one neighbour leaves a nonvanishing monochromatic matching sum
behind, the vertex equation collapses to a single term. That neighbour's column at that
colour then vanishes off the diagonal, and on the diagonal equals the reciprocal of the
surviving sum, so neither is zero.
\end{theorem}

At a site of degree three the assignment of colours to neighbours is a bijection, so every
edge-end receives exactly one colour and the counts match exactly: three site-colour slots
at each of $n$ sites, against $3n/2$ edges with two ends each. Each pinned column carries a
single nonzero entry, on the diagonal --- the local shape of the four-vertex solution,
reproduced at any size.

\begin{remark}[What this does not show]
Sharing that shape on the pinned entries is not being a colour-per-matching system. Such a
system has every entry determined; a spread-case solution agrees with one only where it is
pinned, and its blocks are otherwise unconstrained. So an impossibility theorem about
colour-per-matching does not apply here, and the reduction suggested by the resemblance
does not exist.
\end{remark}

What would close the case is a global obstruction: the two ends of an edge each receive a
colour, and nothing so far forces those to agree or to disagree. Where they agree
throughout, the colour classes are perfect matchings; where they do not, blocks carry two
nonzero diagonal entries instead of one. Both are locally consistent. The obstruction, if
there is one, lives in how the pinned columns interact across sites, and that is where the
argument now stands.

\section{A falsification lane at eight vertices}\label{sec:falsify}

The four-vertex solution is not arbitrary: its vertex set is the Klein group, its
weights are diagonal, and each edge is weighted by a function of the difference of its
endpoints. The natural place to look for a counterexample at eight vertices is the same
construction one dimension up --- the elementary abelian group of order eight --- which
is a family of twenty-one unknowns and can be examined exactly.

Diagonal weights admit only matchings whose edges join equal colours, so a colouring's
amplitude factors over its colour classes, and the conditions become: the matching sum
at each constant colouring is one, and for every splitting into classes of even size the
product of the classes' matching sums vanishes.

\begin{proposition}[Ownership]
Each colour owns a difference: some difference carries nonzero weight at that colour
while every other colour vanishes on it, and the three owned differences are distinct.
\end{proposition}

This follows by expanding the constant-colour condition at a site, which shows some
difference has both nonzero weight and a surviving complementary sum, and then reading
the two-class splittings at that difference.

Three-class splittings say more: two classes of size two contribute a nonzero product,
so the matching sum on the remaining four sites must vanish outright.

Under ownership alone the three-class conditions do not close --- the ideal they
generate together with the normalisation is proper. Solving the full system numerically
over all colourings with even classes, no solution is found: every start collapses to
zero weights, leaving the normalisation unmet. The search was validated first on four
vertices, where a solution exists and is recovered to a residual of order $10^{-30}$, so
the failure at eight vertices is evidence about the family rather than the method.

This is a computation and is recorded as one. It supports a narrow claim worth having: a
counterexample, if one exists, is not the algebraic generalisation of the four-vertex
solution.

\section{Why fewer colours cannot be deformed into more}\label{sec:unused}

Solutions in two colours exist at every even vertex count. It is natural to ask whether
one can be nudged into a three-colour solution, and the answer is that it cannot, for a
reason worth recording.

Read as a three-colour system with the third colour unused, the eight-cycle two-colour
solution satisfies thirteen thousand one hundred and twenty-one of the thirteen thousand
one hundred and twenty-two defining conditions. Every vanishing condition holds, for
every non-constant colouring, and two of the three normalisations hold. Exactly one
equation fails: the monochromatic amplitude at the colour it does not use.

That is not encouragement. The one equation is maximally degenerate there.

\begin{theorem}[An unused colour resists perturbation]
If no edge carries weight at a colour, then perturbing the weights multiplies that
colour's monochromatic amplitude by the perturbation raised to the number of edges in a
matching.
\end{theorem}

So the deficient equation has vanishing gradient, and vanishing derivatives of every
order below the matching size. A path from the two-colour locus cannot reach it; only a
jump can. This was found by attempting the deformation numerically --- a continuation
raising the missing amplitude from zero did not move at the first step, the residual
equalling the target exactly --- and then explained.

The consequence for searching is concrete: the low-colour locus is not a useful starting
point, and a search in a family containing it should begin elsewhere.

\section{Searching the pinned family}\label{sec:pinnedsearch}

The pinning gives a family of candidate weight systems far smaller than the general one,
and defined by a proved statement rather than a guess. Searching it required getting the
formulation right, which took five attempts; the first four each answered a different
question confidently.

Asking directly that the amplitude meet its target fails because the zero weights satisfy
every vanishing condition, so the trivial point is a local minimum and almost every start
falls into it. Removing the scale and asking only that non-constant amplitudes vanish and
the constant ones agree fails because points where \emph{every} amplitude vanishes satisfy
that, and the search finds them. Adding a variable $s$ with $A s = 1$ to force the
amplitude nonzero passes its control but leaves $s$ unbounded, and tracking it shows the
amplitude shrinking as $s$ inflates --- one class appeared to approach a solution with all
three constant amplitudes equal and nonzero, and that was this degeneration. Bounding the
scale alone, without saturation, fails its control outright: on four sites, where a
solution exists, every converged point has vanishing amplitude.

Sphere and saturation together work. The scale is pinned so $s$ cannot inflate, and the
saturation excludes the locus the sphere alone falls into.

\begin{remark}[The separation]
On four sites, where a solution exists, twenty-nine of thirty starts converge and the
largest constant amplitude among them is $0.136$. Across all six cubic classes on eight
sites --- enumerated directly from all $19355$ labelled cubic graphs and reduced by
spectrum --- no start converges at all, in any class, in either the case where the two ends
of each edge agree or in sampled cases where they do not.
\end{remark}

This is a computation and is recorded as one. Its value is the control: four plausible
formulations each returned a confident negative or a spurious near-solution, and each was
caught only by running a case whose answer was known in advance.

\section{What a fully degenerate site forces}\label{sec:fan}

The dichotomy of Section~\ref{sec:site} leaves two branches at each site. The first
asks for a rank drop and is pursued elsewhere. The second says that removing the site
together with any one other site leaves a configuration whose non-constant matching
sums all vanish. Call such a site \emph{fully degenerate}. Read as a disjunct it is
inert; read as a hypothesis it is very strong, and this section extracts what it costs.

Write $\mu_k(x)$ for the matching sum of the constant colouring $k$ on the sites other
than $p$ and $x$. Expanding the amplitude along $p$ shows that these numbers solve a
linear system whose right-hand side is the monochromatic amplitude, so for each colour
there is a partner with a live edge in that colour and $\mu_k \neq 0$ on it.

\begin{theorem}[Purity]\label{thm:purity}
At a fully degenerate site, if $\mu_k(x) \neq 0$ then $W_{p\alpha,\,x\beta} = 0$ for
every $\alpha$ and every $\beta \neq k$.
\end{theorem}

The proof isolates a single term. Colour every site $k$, repaint $p$ with $\alpha$ and
$x$ with $\beta \neq k$; the colouring is non-constant, so its amplitude vanishes.
Expanding along $p$, every partner other than $x$ sees a colouring that is non-constant
on its own complement, which degeneracy annihilates. What survives is one product, and
one of its factors is assumed non-zero.

Two consequences follow at once. A live edge admits only one colour of support, since
purity in two colours leaves the edge nothing to carry; and the three colours therefore
select three \emph{distinct} partners, each joined to $p$ by an edge carrying that
colour alone at the far end. This is precisely the shape of the four-vertex solution,
whose three matchings wear three colours and whose every site meets the other three
along one edge per colour --- and it holds there at all four sites, which has been
checked. The content of the theorem is that no other shape is available, at any size.

\begin{theorem}[Descent]\label{thm:fandescent}
If a fully degenerate site has a partner whose complement retains all three constant
colourings, the complement carries a solution on two fewer sites.
\end{theorem}

At eight sites the complement has six, and there is no six-site solution; so the
descent never fires, and every partner of a degenerate site has deficient support. That
the six-vertex case is settled is exactly what makes this bite.

\subsection*{When every site is degenerate}

The complementary sums are symmetric in their two sites --- the set they run over is
the same one --- so the hypothesis of Theorem~\ref{thm:purity} transports to the other
end of the edge. Purity at one end kills every column but one, purity at the other
kills every row but one, and the two survivors are the same colour.

\begin{theorem}[One colour per edge]\label{thm:onecolour}
If every site is fully degenerate then the block of every live edge which is not inert
has a single entry: there is a colour $\kappa$ with $W_{p\alpha,\,x\beta} = 0$ unless
$\alpha = \beta = \kappa$.
\end{theorem}

Here an edge is \emph{inert} when the sites other than its two ends carry no matching
sum at all; such an edge appears in no amplitude, every term containing it being
multiplied by one of those vanishing sums.

This converts a hypothesis about matching sums into a statement about the weights: the
configuration is a graph with one colour per participating edge, and the amplitude of a
colouring is the product, over the colours, of the matching sums of the corresponding
colour classes on the corresponding parts.

It should be said plainly what this does \emph{not} yet give. One is tempted to argue
that three colour classes, each carrying a perfect matching, must admit a matching using
two colours, and that the associated non-constant colouring would then have non-zero
amplitude. The first half is a combinatorial claim about pairwise edge-disjoint
matchings; the second half is false as stated, because the polychromatic matching is one
term among the compatible ones and the rest may cancel it. So this branch is reduced to
a weighted question about colour classes, not settled by one.

\section{A pure partner, unconditionally}\label{sec:pure}

The results of Section~\ref{sec:fan} pay for their strength with a hypothesis. This
section removes it. The change of method is small and worth stating plainly: instead of
testing the defining equations against a colouring, test them against a \emph{separate
vector at every site}. The vectors are independent of one another, and it is that
independence --- not any assumption about the configuration --- which does the work.

Fix a site $u$ and a colour $\alpha$. Weight each colouring by an indicator pinning $u$
to $\alpha$, and by one test vector $\psi_w$ per remaining site, and sum. Read directly,
the equations collapse the sum: every non-constant colouring contributes nothing, every
constant colouring but one is refused by the indicator, and what remains is
\[
  \mathcal{A}_\alpha \prod_{w \neq u} \psi_w(\alpha).
\]
Read instead through the expansion along $u$, the sum breaks into one term per partner,
and in the $v$-th term the colour at $v$ is summed against $\psi_v$ alone --- the
matching sum on the complement of $\{u,v\}$ does not see it. So the $v$-th term carries
the factor $\langle \psi_v,\, r_v\rangle$, where $r_v$ is the $\alpha$-th row of the
block at $v$, and a test vector annihilating that row kills the term outright.

\begin{theorem}[Pure partner]\label{thm:pure}
For every solution, every site $u$ and every colour $\alpha$ there is a partner $v$ with
$W_{u\alpha,\,vb} = 0$ for all $b \neq \alpha$ and $W_{u\alpha,\,v\alpha} \neq 0$.
\end{theorem}

Suppose not. Then at every partner one may choose a test vector annihilating the row and
having non-zero $\alpha$-component: if the diagonal entry vanishes take the basis vector
itself, and otherwise some off-diagonal entry is non-zero and the two combine. With such
a choice at every partner the expansion gives zero, while the direct reading gives a
product of non-zero numbers. In three-space the only rows admitting no such annihilator
are the non-zero multiples of the $\alpha$-th basis vector, which is the claim.

The hypothesis is exactly the defining equations. Nothing is assumed about supports,
degrees, or the number of sites. The statement is also not vacuous: it fails on random
weights, and it holds at every site and colour of the four-vertex solution, where the
partner it names is the one joined by that solution's edge of that colour.

Two limits should be stated. The theorem produces \emph{one} partner per site and
colour, not all of them, and it constrains a row rather than a block --- the partner it
names need not be pure when read from the other end. Under the degeneracy hypothesis of
Section~\ref{sec:fan} both limits lift, giving three distinct partners and single-entry
blocks. Whether they lift unconditionally is open; on four sites every solution the
search returns has single-entry blocks throughout, which is evidence about four sites and
about the formulation searched, not a theorem.

\section{The master relation}\label{sec:master}

Section~\ref{sec:pure} tested the equations at a site against a basis vector. This
section tests them against an arbitrary covector, and the gain is not incremental: a
basis vector isolates one colour, a general covector couples all three at once, and it
is that coupling which produces a three-fold conclusion.

Fix a site $u$ and a covector $\rho$. Weight each colouring by $\rho$ at $u$ and by a
test vector $\psi_w$ at every other site, and sum. Read directly, only the constant
colourings survive and the sum is $\sum_k \rho_k \mathcal{A}_k \prod_{w \neq u}
\psi_w(k)$. Read through the expansion along $u$, it becomes one term per partner, and
the $v$-th term carries the factor $\langle \psi_v,\, \rho^{\!\top} W_{uv}\rangle$: the
colour at $v$ meets its own test vector and nothing else, because the matching sum on the
complement of $\{u,v\}$ never sees it.

\begin{theorem}[Master relation]\label{thm:master}
If $\langle \psi_v,\, \rho^{\!\top} W_{uv}\rangle = 0$ for every $v \neq u$, then
\[
  \sum_{k} \rho_k\, \mathcal{A}_k \prod_{w \neq u} \psi_w(k) \;=\; 0 .
\]
\end{theorem}

Write $r_w := \rho^{\!\top}W_{uw}$, a covector in three-space, and $L_w := \ker r_w$, of
dimension at least two. The relation holds for every choice of $\psi_w \in L_w$.

Two consequences are immediate. Testing every site with a colour's indicator gives the
\emph{column lemma}: a covector annihilating one and the same colour's column at every
partner has no component in that colour. And the whole of Section~\ref{sec:pure} is the
case $\rho = e_\alpha$.

\subsection*{The fan for a generic covector}

\begin{theorem}[Generic fan]\label{thm:rhofan}
Let $\rho$ have no vanishing component, and let the site have at least two partners. Then
for every colour $k$ there is a partner $w$ with $r_w$ a non-zero multiple of $e_k$; and
the three partners so named, one per colour, are distinct.
\end{theorem}

Suppose no partner qualified for the colour $k$. Then at each $w$ one may choose
$\psi_w \in L_w$ with $\psi_w(k) \neq 0$: if $r_w(k) = 0$ take $e_k$ itself, and otherwise
some other coordinate of $r_w$ is non-zero and the two combine. That alone is not enough
--- the relation still has three live terms. What finishes it is that the other two terms
can be killed at two \emph{different} partners.

Write $k'$ and $k''$ for the other two colours. If some partner has $r_w(k) = 0$, then
$e_k$ itself lies in $L_w$ and kills both other terms at once. Otherwise every partner has
$r_w(k) \neq 0$, and then no partner can have both $r_w(k')$ and $r_w(k'')$ zero --- that
would make $r_w$ a non-zero multiple of $e_k$, which was excluded. If $r_w(k'') = 0$ at
every partner the column lemma gives $\rho_{k''} = 0$, against the hypothesis; likewise
for $k'$. So some partner has $r_w(k'') \neq 0$ and some partner has $r_w(k') \neq 0$, and
if these coincide, a third partner --- available because the site has at least two ---
supplies the other. At a partner with $r_w(k'') \neq 0$ the vector
$r_w(k'')e_k - r_w(k)e_{k''}$ lies in $L_w$, has non-zero $k$-component and vanishes in
the colour $k'$; symmetrically for $k''$. With these choices the relation reads
$\rho_k \mathcal{A}_k \prod_w \psi_w(k) = 0$, a product of non-zero numbers.

Distinctness is immediate: a covector cannot be a non-zero multiple of two different basis
vectors.

\subsection*{From covectors to blocks}

The partners named by Theorem~\ref{thm:rhofan} depend on $\rho$. They can be made not to.

For a partner $w$ and a colour $b \neq k$, the covectors annihilating that column form
the kernel of a linear form; a partner fails to be named only if one of those forms is
non-zero, and then its kernel is a proper subspace. Finitely many proper subspaces cannot
cover a space over an infinite field, so a covector can be chosen off all of them and off
the three coordinate hyperplanes at once. Theorem~\ref{thm:rhofan} applied to that
covector names a partner, and the forms it annihilates must therefore have been zero
outright --- which is a statement about the block, with no covector left in it.

\begin{corollary}[Column-supported blocks]\label{cor:colblock}
For every site $u$ and every colour $k$ there is a partner $w$ whose block satisfies
$W_{ua,\,wb} = 0$ for every $a$ and every $b \neq k$, with the block not identically zero.
The three partners so named are distinct.
\end{corollary}

Read from the far end this says the block $W_{wu}$ is supported in a single row. The
four-vertex solution realises it exactly: each of its three edges at a site carries one
colour, and the three partners are the three other sites.

\subsection*{Cancellation, turned around}

The remaining steps must confront cancellation directly, and there is a way to do so that
does not estimate anything.

\begin{lemma}[No rival, no cancellation]\label{lem:unique}
If a colouring admits exactly one matching with non-zero contribution, that colouring is
constant.
\end{lemma}

The proof is one line --- the amplitude equals that contribution --- but the use is not.
It converts a combinatorial statement, that some matching is alone in being compatible
with a colouring, into an algebraic one about the equations. That is the intended route
for what remains.

\subsection*{What remains}

Two steps stand between Corollary~\ref{cor:colblock} and the theorem. The second is
settled here; the first is not.

\paragraph{The steps that are settled.} Two of the three remaining links are done.

\begin{lemma}[A fourth matching]\label{lem:fourth}
A cubic graph on $2n \geq 6$ vertices which is the union of three disjoint perfect
matchings has a perfect matching other than those three.
\end{lemma}

Write the classes $M_0, M_1, M_2$. If $M_0 \cup M_1$ has two or more cycles, exchanging
the colours along one cycle produces a fourth matching, so it is a Hamiltonian cycle $C$;
label its vertices $1,\dots,2n$ so that $M_0$ joins $2i-1$ to $2i$.

Each edge of $M_2$ is a chord. If a chord joins two vertices at odd distance along $C$,
deleting its ends leaves two paths with evenly many vertices, each of which its own
alternating edges match; together with the chord this is a fourth matching. So every
chord joins vertices of equal parity, and $M_2$ splits into a perfect matching of the odd
positions and one of the even positions.

If an odd chord and an even chord \emph{interleave} along $C$, deleting all four ends
leaves four paths, each with evenly many vertices --- the ends of each chord agree in
parity and the two chords disagree, so every gap is even. Matching the paths alternately,
with the two chords, gives a fourth matching. So no odd chord interleaves an even chord.

That is impossible, and a single comparison shows it. Take the coordinates the two classes
supply: $x_i$ joined to $y_i$ by the first matching and to $y_{i+1}$ by the second, indices
running over $\mathbb{Z}/m$ with $2m$ the number of vertices. Among all chords of the
$x$-class choose one of shortest forward span $d$, from $x_{i_0}$ to $x_{i_0+d}$, and stand
at $w = i_0 + 1$, the index just past its tail. The $y$-chord at $w$ runs forward some
distance $v \geq 1$, to $y_{w+v}$.

If $v < d$, that chord stops strictly inside the span, and the $x$-chord \emph{at $w$}
interleaves with it: its near end is $x_w$, inside, and its far end lies at forward
distance at least $d > v$, outside, because $d$ was the shortest span of any $x$-chord. If
instead $v \geq d$, the $y$-chord reaches past $x_{i_0+d}$, which sits at forward distance
$d-1$ from $w$, but not past $x_{i_0}$, a full turn away at distance $m-1 \geq v$; so the
shortest chord itself interleaves with it. Either way an $x$-chord and a $y$-chord
interleave, against the assumption. Minimality is spent once, on a single lower bound.

This whole proof is machine-checked, coordinates and all. Both matchings --- the one built
from a chord across the classes and the one built from two interleaved chords --- are
constructed explicitly and verified compatible and non-constant; the comparison above is
formalised as the statement that interleaved chords always exist; and the coordinates
themselves are built from an abstract pair of disjoint involutions with connected union, as
the orbit of their product together with its image under the first. The two classes are
disjoint because a site of the orbit landing back inside it would force one of the two
involutions to fix a point: the doubling map on the orbit hits one of any two consecutive
residues. Connectivity makes the union everything.

\begin{theorem}[The rigid shape is refuted]\label{thm:mixed}
Suppose every live edge joins a site to its partner in one of three disjoint perfect
matchings and carries that matching's colour at both ends. Then no perfect matching mixes
two of the three.
\end{theorem}

This is machine-checked. A mixed matching induces a colouring which is non-constant, and
for which it is the \emph{only} contributing matching: any other sends some site to a
partner whose edge carries a different colour, and that entry is zero. The no-rival lemma
then makes the colouring constant, a contradiction. With Lemma~\ref{lem:fourth} this says
the rigid shape cannot be sustained above four sites.

\paragraph{Where the argument now stands.} Call a configuration \emph{three-regular} when
every site has exactly three live partners, one column-supported in each colour. The
following is machine-checked, from the defining equations down.

\begin{theorem}[The rigid shape, and its refutation]\label{thm:threereg}
A three-regular solution is three disjoint perfect matchings whose edges are the only live
ones, each carrying its own colour at both ends; and no perfect matching mixes two of the
three.
\end{theorem}

Diagonality is the step that had looked hardest and turns out to be free. If a site's only
live partners are its three column-supported ones, a covector orthogonal to one block's
live column is orthogonal to \emph{every} block's column in that colour --- the other two
blocks have none there --- and the column lemma then says the covector has no component in
that colour. Two hyperplanes containing each other are equal, so the column is a multiple
of the colour's basis vector: the block is a single entry on the diagonal.

From there the shape assembles. A single diagonal entry reads the same from either end, so
a site's colour-$k$ partner has it back as its own colour-$k$ partner; the partner maps are
fixed-point-free involutions, that is, perfect matchings; they are disjoint because the
three partners are distinct; and every live edge is one of theirs. A mixed matching is then
alone on its colouring, which the equations refuse.

So the whole conjecture rests on exactly two facts.

\begin{enumerate}
\item A solution is three-regular. \emph{This is the gap.} The column corollary supplies
at least three column-supported partners; what is missing is that there are no others ---
neither a partner outside them nor a second one serving the same colour. Both cases are
one statement: for such a partner $v$, every colour still has a column-supported partner in
$V \setminus \{u,v\}$, and the spared form of the master relation then forces the weighted
matching sum of the complement to vanish identically on the product of the kernels $L_w$.
Excluding that vanishing is all that remains.
\item A three-regular configuration on six or more sites has a mixed matching. This is
Lemma~\ref{lem:fourth}, proved and formalised above.
\end{enumerate}

\subsection*{Inside the gap}

The gap splits, and one half yields.

\begin{lemma}[Orthogonality localises the witness]\label{lem:localise}
If a covector annihilates one colour's column at every partner outside a set, and has a
non-zero component in that colour, then some partner inside the set has its contracted row
pointing along that colour.
\end{lemma}

This is machine-checked, and the column lemma is its empty-set case. Testing every partner
outside the set with the colour's basis vector is legitimate exactly under the
orthogonality hypothesis, and it kills the other two colours outright --- each of their
products picks up a zero factor from any site outside the set.

Take the set to be the partners in question and the colour $\alpha$. The named partners of
the other two colours have no $\alpha$-column, so the hypothesis reduces to
$\rho \perp c_\alpha$, where $c_\alpha$ is the live column of the $\alpha$-named block.
Hence: for every $\rho$ orthogonal to $c_\alpha$ with $\rho_\alpha \neq 0$, one of the
partners in question is an $\alpha$-witness. Covering the hyperplane
$\{\rho \cdot c_\alpha = 0\}$ by the subspaces of covectors naming each, one partner names
for the whole hyperplane, so its row space carries $\alpha$'s axis on it.

If that happens in all three colours at the \emph{same} partner, its row space lies in
$\langle e_0,e_1\rangle \cap \langle e_0,e_2\rangle \cap \langle e_1,e_2\rangle = 0$, so
its block vanishes and it was never live. In particular, when there is only one partner
in question, the gap closes.

What the lemma cannot see is the case $c_\alpha \parallel e_\alpha$ --- the named block
already a single diagonal entry. Then $\rho \perp c_\alpha$ forces $\rho_\alpha = 0$ and
the hypothesis is unsatisfiable. So the surviving case is: named blocks already diagonal,
with extra live partners beside them. There the relation at the site is exhausted; every
choice of test vectors is annihilated by one of the three named partners, and no
information is left in it.

One reduction applies there. If the complement of a pair carries no matching sum at all,
the edge appears in no amplitude and may be set to zero without changing any of them; such
edges can be removed. So the surviving case may be taken to have every live edge
participating somewhere --- which is what makes the vanishing that defines the gap a
statement with content rather than a degeneracy.

Two further routes into the surviving case were followed and are recorded as closed.
Sparing \emph{two} named partners at once, with the covector taken to be a colour's basis
vector, does give a clean relation: the complementary sum vanishes unless the other named
partners wear that colour. But it is the rank-one determination already in hand, read
coefficient by coefficient, not a new equation --- and its content, that the complement of
a site and its colour-$\alpha$ partner is supported on the constant colouring $\alpha$,
is a consequence of the shape rather than a step towards it. Pairing the relation at a
site with the relation at a free partner is likewise circular, whereas pairing it with the
relation at a \emph{named} partner is not, and gives the dichotomy recorded above.

\subsection*{The gap as a tensor membership}

The right frame for what remains is not a further specialisation of the relation but a
statement about the complementary matching sum as a \emph{tensor}.

Read $H$, the amplitude tensor of the complement of $\{u,v\}$, as a multilinear form on
one three-space per remaining site. The gap says $H$ vanishes on $\prod_w L_w$, a product
of hyperplanes $L_w = \ker \ell_w$ with $\ell_w = \rho^{\!\top}W_{uw}$. Splitting each
three-space as $L_w \oplus \langle e_w\rangle$ with $\ell_w(e_w) = 1$ and expanding shows
that this is exactly membership in the ideal those forms generate:
\[
  H\big|_{\prod_w L_w} = 0
  \quad\Longleftrightarrow\quad
  H = \sum_{w} \ell_w \otimes G_w ,
\]
with $G_w$ a tensor on the remaining sites. What has to be shown is that a \emph{matching}
tensor admits no such decomposition --- that its components $\operatorname{pm}(d)$ cannot
be written as $\sum_w \ell_w(d_w)\,(G_w)_{d}$.

Three consequences of the membership are already visible, and each is a special case of
contracting all but a few sites with kernel vectors. Contracting all of them gives the gap
itself. Contracting all but one gives a functional at the spared site proportional to
$\ell_w$ --- so at a named partner, whose contracted row points along $e_\beta$, the
combination $\sum_z (W_{\beta z}\psi_z)J_z$ points along $e_\beta$ too; and at a site dead
for $u$, where $\ell$ vanishes, that combination vanishes outright. Contracting all but two
gives a sum of two rank-one tensors with prescribed factors, which is the dichotomy
recorded above.

The decomposition has one immediate and useful feature: at a site \emph{dead} for $u$ the
form $\ell_w$ vanishes, so that term is absent. The sum runs only over $u$'s live partners
in the complement --- in the surviving case, the three named ones and whatever extra
partners remain.

\begin{corollary}[Support]\label{cor:support}
If $u$ has exactly one extra live partner $v$, then in the complement of $\{u,v\}$
\[
  \operatorname{pm}(d) = 0
  \qquad\text{whenever } d(w_\beta) \neq \beta \text{ for all three } \beta .
\]
\end{corollary}

Only the three named terms survive, and in the surviving case $\ell_{w_\beta}$ is a
multiple of $e_\beta$, so each term carries the indicator of $d(w_\beta) = \beta$. The
statement itself no longer mentions the covector.

What it costs to contradict is worth stating exactly, because the natural reading
overstates it. Take any perfect matching $N$ of the complement built from live edges and
colour each site by its $N$-edge. Named edges compatible with that colouring are exactly
$N$'s own, since a site wearing $\beta$ has only one colour-$\beta$ partner and it is the
one $N$ gave it. Each $w_\beta$ is then automatically coloured off its own colour --- its
colour-$\beta$ partner is $u$, which is absent --- \emph{provided} $N$ matched it by a
named edge. So if $N$ were the only contributing matching, its product would be non-zero
and the corollary contradicted.

Uniqueness is where it stops. Extra edges elsewhere in the complement may also be
compatible with that colouring, and nothing so far excludes them. The argument therefore
closes only when the complement has no extra edges at all --- when $uv$ is the sole extra
edge of the whole configuration, in which case either such an $N$ exists and the
corollary is contradicted, or none does and every matching sum of the complement vanishes,
making $uv$ inert and removable. Either way that configuration is refuted.

The support statement is in fact sharper than the covector reading gives, and comes out
of the spared relation by testing against \emph{basis} vectors instead of general ones.
Each partner is then pinned to a colour, the annihilation condition becomes the vanishing
of a single weight, and the complementary sum collapses from a weighted average to a
single matching sum.

\begin{theorem}[Support]\label{thm:support}
Let $\alpha$ be a colour whose row from $u$ toward $v$ does not vanish. If a colouring of
the complement of $\{u,v\}$ places every partner at a colour where that row vanishes, and
is not constant $\alpha$, then its matching sum on the complement is zero.
\end{theorem}

This is machine-checked. It is sharper than before because only the colour-$\alpha$ named
partner is constrained: the other two named blocks are single entries in their own
colours, so they have no $\alpha$-row at all and impose nothing.

Four consequences follow in the surviving case, and together they say a good deal about
what an extra edge would have to look like.

First, an extra block has at least two non-zero rows as well as at least two non-zero
columns. Being an extra partner is symmetric here: if $u$ were a named partner of $v$ its
block would be a single entry, hence column-supported from $u$ as well, and $v$ would not
be extra.

Second, therefore, for at least two colours $\alpha$ the theorem applies, and a colouring
of the complement with non-vanishing matching sum must place $w_\alpha$ at $\alpha$ for
each of them.

Third, such a $w_\alpha$ must be matched by an \emph{extra} edge. Its colour-$\alpha$
partner is $u$, which is absent from the complement, and its other named edges require it
to wear a different colour.

Fourth, if some such $w_\alpha$ has no extra partner in the complement, every matching sum
there vanishes, the edge $uv$ appears in no amplitude, and it may be deleted. So in a
configuration with no removable edge, every site carries an extra partner --- the property
propagates along the named matchings, which connect the sites.

The same theorem also shows that when $u$ has just one extra partner, that partner
contributes nothing to any constant colouring: taking $\alpha$ different from the constant
puts every named partner at a colour the $\alpha$-row misses.

So an extra edge is far from free.

\subsection*{The gap as a determinantal condition}

The covector in the gap hypothesis is existentially, not universally, quantified, and
saying so exactly turns the condition into a determinantal one.

Fix $u$, an extra partner $v$, and test vectors $\psi$ at the remaining sites. Put
$z_w := W_{uw}\psi_w \in \mathbb{C}^3$, so that the annihilation condition
$\psi_w \perp \rho^{\!\top}W_{uw}$ reads $\rho \perp z_w$. The spared relation is
\[
  \bigl(\chi^{\!\top}(\rho^{\!\top}W_{uv})\bigr)\, H(\psi)
  \;=\; \sum_k \rho_k \mathcal{A}_k\, \chi_k \prod_{w} \psi_w(k) ,
\]
so $H(\psi) = 0$ follows from two things: that $\prod_w \psi_w(k)$ vanishes for every
colour --- a condition on $\psi$ alone --- and that \emph{some} $\rho$ orthogonal to all
the $z_w$ has $\rho^{\!\top}W_{uv} \neq 0$.

In the surviving case both simplify sharply. Taking $\psi_{w_\beta}(\beta) = 0$ at each
named partner secures the first condition, and it makes $z_{w_\beta} = 0$: the named block
is $c_\beta \otimes e_\beta^{*}$, so $z_{w_\beta} = \psi_{w_\beta}(\beta)\,c_\beta$. Sites
dead for $u$ give $z_w = 0$ outright. So only the \emph{other extra partners} contribute
to the span, and the second condition is exactly
\[
  \operatorname{im} W_{uv} \;\not\subseteq\; \operatorname{span}\{\,z_{v'} : v' \in F(u)
  \setminus \{v\}\,\} ,
\]
the image being the column space, which lives in the same three-space as the $z$'s --- the
one indexed by $u$'s colour. (A covector kills $\rho^{\!\top}W_{uv}$ exactly when it
annihilates that column space, not the row space.) When $W_{uv}$ has rank three the
condition is precisely $\dim \operatorname{span}\{z_{v'}\} \leq 2$: the vanishing of every
$3 \times 3$ minor of the $z$'s.

\begin{theorem}[Span domination]\label{thm:spandom}
In the surviving case, if the complementary sum does not vanish then
$\operatorname{im} W_{uv}$ lies in the span of the columns the test vectors pick out at
the remaining partners.
\end{theorem}

This is machine-checked, and the proof is separation rather than genericity: a column
outside that span admits a functional annihilating the span but not it, and the covector
read off that functional empties the relation. No determinantal ideal and no radicality
enter. The three named partners cost nothing --- annihilation at the colour-$k$ named
partner reads $\rho_k \, W_{uk,\,w_k k}\, \psi_{w_k}(k) = 0$, so either the covector
vanishes in that colour or the test vector does, and the $k$-th direct term carries
whichever factor vanishes. And once the test vector at a named partner sits off its own
colour, the column it picks out is zero outright, so the span is carried by the extra
partners alone.

On basis test vectors this becomes a statement with no covector and no test vectors in
it at all. A colouring $c$ of the complement selects one column at each partner, and the
complementary sum collapses to the single matching sum $\operatorname{pm}(c)$:

\begin{corollary}[Span domination, on colourings]\label{cor:spancol}
If $\operatorname{pm}(c) \neq 0$ on the complement of $\{u,v\}$, then
\[
  \operatorname{im} W_{uv} \;\subseteq\;
  \operatorname{span}\bigl(\{\,e_k : c(w_k) = k\,\} \cup
    \{\,\text{the } c(v')\text{-column of } W_{uv'} : v' \in F(u)\setminus\{v\}\,\}\bigr).
\]
\end{corollary}

The named partners contribute $e_k$ when the colouring puts them \emph{on} their own
colour and nothing otherwise, because their blocks are single diagonal entries. So a
colouring placing every named partner off its colour leaves only the extra partners'
selected columns --- and if $u$ has no other extra partner, leaves nothing, forcing
$W_{uv} = 0$. That is the support theorem, recovered; and the general statement bounds
the rank of the spared block by the number of other extras.

Read against a coordinate functional it sharpens once more. Every column the colouring
selects misses the $a$-coordinate unless the $a$-named partner sits on colour $a$; so a
live entry in row $a$ of the spared block forces exactly that.

\begin{theorem}[The live rows are pinned]\label{thm:pinned}
Suppose the three named blocks at $u$ are single diagonal entries and $u$'s only extra
partner is $v$. Then every colouring of the complement with non-vanishing matching sum
places the named partner of each live row of $W_{uv}$ on that very colour.
\end{theorem}

The diagonality hypothesis is not decoration: span domination as proved uses it, to turn
annihilation at a named partner into the vanishing of that colour's coefficient. So this
theorem cannot be used to \emph{derive} diagonality --- a point on which an earlier draft of
this section was circular, and which the next paragraph repairs.

This is machine-checked. An extra block has at least two live rows --- it is
column-supported from neither end --- so every surviving colouring pins at least two named
partners at once. And a named partner pinned to its own colour must be matched by an extra
edge, since the partner it would otherwise use in that colour is $u$, which is absent.

So the gap has become: exhibit one colouring of the complement, with non-vanishing matching
sum, that leaves some live row's named partner off its colour. Equivalently, show the edge
is inert, and delete it. Nothing analytic remains --- what blocks it is that a
non-vanishing matching sum still has to be certified, and the only certificate available
is a compatible matching with no rival.

\paragraph{Tracing non-vanishing back to the constants.} The only non-vanishing the
equations hand us sits at the constant colourings, so every certificate has to be traced
back there. A colouring differing from a constant at exactly one site can be: expand at the
flipped site, and its edge carries the sole off-colour weight with a monochromatic sum
behind it. This is machine-checked, and so is its consequence.

\begin{theorem}[Flip orthogonality]\label{thm:flip}
Fix a site $w$ and a colour $k$. Let $m_z$ be the colour-$k$ matching sum on the complement
of $\{w,z\}$. Then $\sum_z W_{(w,a)(z,k)}\, m_z$ vanishes for every $a \neq k$ and equals
the colour-$k$ amplitude, hence is non-zero, for $a = k$.
\end{theorem}

Column support confines that sum. The named partners of the other two colours contribute
nothing, so when a site has a single extra partner $y$ the relation has just two terms:
\[
  W_{(w,a)(w_k,k)}\, m_{w_k} \;+\; W_{(w,a)(y,k)}\, m_y \;=\; A_k\,\delta_{a,k}.
\]
Read it at $m_y = 0$. The extra drops out entirely; the diagonal case $a = k$ makes
$m_{w_k} \neq 0$, and then every off-diagonal case forces $W_{(w,a)(w_k,k)} = 0$.

\begin{corollary}[Diagonality from a vanishing complementary sum]\label{cor:diagzero}
If a site has exactly one extra partner and the colour-$k$ matching sum on the complement of
the site and that extra vanishes, then the colour-$k$ named block is the single diagonal
entry.
\end{corollary}

This is machine-checked, and it uses nothing but the equations, the column fan and the flip
relation --- no diagonality assumed anywhere, so nothing is circular. Diagonality used to be
available only downstream of full three-regularity; it is now available one colour at a
time, from a vanishing complementary sum.

\paragraph{The spared relation for a covector.} The obstacle to going further was that the
spared relation existed only one row at a time, and a single row cannot annihilate a
subspace --- which is exactly why span domination had to assume diagonality. The rows do
combine. Expanding at $u$ against a covector $\rho$,
\[
  \sum_{y \neq u} \langle \rho, W_{uy}\psi_y\rangle\, K_y(\psi)
    \;=\; \sum_a \rho_a A_a \prod_{w \neq u} \psi_w(a),
\]
where $K_y$ is the matching sum on the complement of $\{u,y\}$ read against the test
vectors. The one obstruction to adding the rows is bookkeeping: the paired sum carries a
colour pin at each end of the spared edge. A swap at each end shows it does not depend on
them. This is machine-checked, and so are its two readings.

\begin{theorem}[Span domination, unconditional]\label{thm:spandom}
If a colouring of the complement misses every colour somewhere and its matching sum does not
vanish, the column image of the spared block lies in the span of the columns that colouring
selects. Nothing is assumed about the named blocks.
\end{theorem}

\begin{theorem}[The constant reading]\label{thm:constread}
For a colour $k$ and a covector $\rho$ annihilating every complement site's colour-$k$
column, $\langle \rho, W_{uv}x\rangle$ times the colour-$k$ matching sum on the complement
equals $\rho_k A_k x_k$.
\end{theorem}

With column support the annihilation hypothesis is cheap. At a site whose live partners are
its three named ones and a single extra $v$, only the colour's own named partner has a
non-zero colour-$k$ column, namely $v_k = W_{(u,\cdot)(w_k,k)}$; so any $\rho$ orthogonal to
that single vector qualifies. This is what the count of degenerate colours needed.

\begin{proposition}[Some colour is degenerate]\label{prop:onedegen}
At a site whose live partners are its three named ones and a single live extra, the
complementary sum at the extra vanishes for at least one colour --- and that colour's named
block is therefore the single diagonal entry.
\end{proposition}

Suppose instead $m_v^k \neq 0$ for all three colours. Take $x$ with $x_k = 0$: the constant
reading gives $\langle \rho, W_{uv}x\rangle = 0$ for every $\rho \perp v_k$, so each column
$c_b$ of the extra block with $b \neq k$ lies on the line $\mathbb{C}v_k$. The flip relation
$A_k e_k = m^k_{w_k} v_k + m^k_v c_k$ puts $e_k$ in $\mathrm{span}\{v_k, c_k\}$, hence in
$\mathrm{span}\{v_k, v_{k'}\}$ for either other colour. If the three $v_k$ lay in a plane,
all three $e_k$ would too; they do not, so the $v_k$ are a basis. But then each column $c_b$
lies on both of the two distinct lines $\mathbb{C}v_{k'}$, $\mathbb{C}v_{k''}$ with
$k',k'' \neq b$, so $c_b = 0$. The extra block vanishes, against the extra being live.

This is machine-checked, together with the two theorems and the flip relation it uses.

\paragraph{Minimise the support first.} The degree bound should not be attacked over every
presentation of a solution. Any hypothetical solution has a support-minimal representative on
the same vertex set --- zeroing a block strictly lowers the total live degree, and the process
must stop --- and in that representative every live pair is provably non-inert, because an
inert block could be zeroed without changing any amplitude. So the certificate the
support-theoretic arguments keep needing is not an extra assumption; it is available on every
live edge for free. The conjecture follows from the degree bound imposed only on
support-minimal representatives, and that is where the remaining attack belongs.

Two consequences are already in hand. A support-minimal system has no bridge: a cut cannot
have a single live crossing pair, since minimality makes that pair participate, which forces
the cut to be odd, and then every matching uses the pair --- three non-zero monochromatic
amplitudes demand three diagonal colours on it, while a participating complementary colouring
permits only one.

And the site-local machinery now consumes the certificate directly.

\begin{theorem}[Two terms, and a transfer]\label{thm:split}
At a site whose live partners are its three named ones and a single extra, a colouring
sending the site and its $\alpha$-named partner both to $\alpha$ splits the amplitude into
exactly two pieces: the named edge times its complement's matching sum, and the extra's edge
times its complement's. On a non-constant colouring the amplitude vanishes, so the two pieces
cancel. A dead entry at the colouring kills the named partner's complement; a live entry
against a participating extra hands the certificate on to the named pair.
\end{theorem}

Machine-checked. Diagonality kills every other named partner on the row and every other site
has no live edge, which is why the split is exactly two terms and not four.

The same expansion, run at a colouring that sends the site to $\alpha$ but its
$\alpha$-named partner \emph{elsewhere}, kills the named term as well: the other two die on
the row, and this one because the partner is off its colour. Only the extra survives, so a
non-vanishing complementary sum there forces that row entry to zero.

\begin{theorem}[Pinning from a certificate]\label{thm:certpin}
On a colouring whose complementary sum at the extra does not vanish, every live row of the
extra block pins its named partner to that row's colour.
\end{theorem}

Machine-checked, and the derivation is now three lines of expansion --- no separation, no span
domination, and therefore none of the diagonality hypothesis that made the earlier route
circular. Combined with minimality, the statement has real force: on a minimal representative
the certificate exists at every live pair, so the pinning applies unconditionally there. In
the case where every colour is degenerate at a site --- where the constant colourings supply
nothing --- minimality still produces a non-constant certificate, and every live row of the
extra block is pinned by it. The named partner's complementary sum is then determined
outright: it is the extra's row entry at that colouring, times the extra's complementary sum,
over the named weight.

\paragraph{Degree propagates across a thick edge.} The two ends of an edge constrain each
other, and the constraint is exactly about degree.

\begin{theorem}[Both ends]\label{thm:bothends}
If the block over a live pair has two distinct live rows, then the far end has at least four
live neighbours.
\end{theorem}

Machine-checked. Neither endpoint can see the other as a column-supported partner when two
rows are live, so the column fan at the far end names three partners, none of them this one.
The contrapositive is the useful reading: at a site of degree exactly three, every live block
is confined to a single row when read from the far end.

So a fourth live partner is not a local accident. Either the extra block is \emph{thin} ---
one live row, and then the far end may well be three-regular --- or it is thick, and the
excess degree propagates across the edge. That splits the hard case in two, and it suggests
looking for the contradiction along the propagation rather than at a single site.

Applying the bound at both ends at once settles the three-regular part completely.

\begin{corollary}[Two three-regular ends]\label{cor:singleentry}
A block between two sites of degree three is a single entry.
\end{corollary}

Machine-checked. Degree three at the far end confines the block to one row, degree three at
the near end to one column, and one row meeting one column leaves one entry. So no colour
mixing can hide in the three-regular part of a configuration: every edge there is
monochromatic at each end, and whatever is left to rule out lives entirely on the sites of
excess degree.

And in the hard case the excess sits at \emph{both} ends, whether the block is thick or thin.

\begin{theorem}[A three-regular end forbids degeneracy]\label{thm:threeregend}
At a site of degree three the live partners are exactly the three the fan names, so its
blocks are single diagonal entries. The flip relation there, in the colour naming a given
edge, has one surviving term --- the edge back, times the complementary sum --- and it equals
the amplitude. So that complementary sum is non-zero.
\end{theorem}

\begin{corollary}[Excess at both ends]\label{cor:bothexcess}
A site whose complementary sums with an extra vanish in every colour cannot have that extra
be three-regular.
\end{corollary}

Machine-checked. The earlier propagation covered thick blocks only; this covers thin ones as
well, and it does so through the far end's own equations rather than through the fan count.
In the remaining configuration, then, both endpoints of the extra edge carry excess degree.

\paragraph{Excess degree cannot sit alone.} One three-regular end already makes a block a
single diagonal entry, so every edge \emph{touching} a three-regular site is monochromatic in
one colour at both ends --- not merely the edges between two of them. That has a sharp
consequence.

\begin{theorem}[No isolated excess]\label{thm:noisolated}
A site with four live partners has a live partner that also has four.
\end{theorem}

Machine-checked. Suppose not. Then every block at the site is a single diagonal entry, and
four blocks in three colours must repeat a colour. The two partners sharing that colour are
each left with this site as their \emph{only} live partner in it, because a three-regular
site's colour classes are singletons. No perfect matching in that colour can match both to
the same site, so every term of that colour's amplitude vanishes --- and the monochromatic
amplitudes are not zero.

The mechanism deserves its own statement, because it is used three times over: two distinct
sites cannot share their unique live partner in one colour, since a perfect matching in that
colour would have to send both to it. From that the count is immediate.

\begin{theorem}[At most three]\label{thm:atmostthree}
A site's three-regular live partners occupy distinct colours, so it has at most three of them
--- and a site of degree $d$ has at least $d-3$ partners of excess degree.
\end{theorem}

Machine-checked. Globally the same lemma says: for each colour, distinct three-regular sites
have distinct partners in that colour, so the colour's perfect matching is forced on the whole
three-regular part at once.

A naive count does not close from here, and it is worth recording why. Writing $T$ for the
three-regular sites and $D$ for the rest, the three edges at each $T$-site give
$3|T| = 2e(T) + e(T,D)$, the bound above gives $e(T,D) \leq 3|D|$, and the degree bound gives
$\sum_{d \in D}(\deg d - 3) \leq 2e(D)$. These are mutually consistent for every split, so the
contradiction has to come from the colour structure inside $D$, not from edge counting.

The sites of excess degree therefore have no isolated member, and the search for the
contradiction can be confined to the subgraph they span.

The same rigidity closes another escape, this time inside the hard case itself.

\begin{theorem}[A three-regular named partner destroys the certificate]\label{thm:namedkill}
If the named blocks at a site are single diagonal entries, a colouring of the complement
sends the $\alpha$-named partner to colour $\alpha$, and that partner is three-regular, then
the complement's matching sum vanishes.
\end{theorem}

Machine-checked. A three-regular site's only live partner in colour $\alpha$ is the one its
own block names, and here that is $u$ --- which the complement excludes. The partner can be
matched to nothing, so every term dies.

Read against the pinning, this says that in the hard case every named partner of a live row
of the extra block carries excess degree. Together with the extra itself, whose degree is
already known to exceed three, the site has at least two excess neighbours. The excess
subgraph is not a tree of pendant attachments; every site in it that carries a degenerate
extra has at least two neighbours inside it.

\paragraph{No set is closed in two colours.} The obstruction the programme has been looking
for turns out to be a statement about cuts rather than cycles.

\begin{theorem}[Two-colour closure]\label{thm:twocolour}
Let $C$ be a non-empty proper set of sites with no live colour-$k$ edge and no live
colour-$l$ edge leaving it, and no live edge joining colour $l$ inside to colour $k$ outside.
Then the equations are contradictory.
\end{theorem}

Machine-checked. Both monochromatic amplitudes factor across the cut, so all four factors are
non-zero. The colouring wearing $l$ inside and $k$ outside respects the same cut and is not
constant, so its amplitude vanishes --- while factoring into two of those very factors.

This is the alternating cycle, in set form. A cycle alternating two colours inside the
three-regular part is exactly such a set: each of its sites has a unique partner in each of
the two colours, both on the cycle, and its blocks are single diagonal entries so no mixed
edge leaves. So no such cycle exists.

The consequence for the three-regular part is immediate. For each pair of colours, the graph
of internal edges of those two colours has maximum degree two and no cycle, hence is a union
of paths. Writing $a_k$ for the number of three-regular sites whose colour-$k$ partner lies
outside the three-regular part, the number of paths is $(a_k+a_l)/2$, so $a_k + a_l \geq 2$
for every pair. The three-regular part cannot be closed under any two colours at once.

\paragraph{The named partners cannot all be three-regular.} The rigidity results combine into
a closure.

\begin{theorem}[Not all three named]\label{thm:notallnamed}
A site with exactly one extra partner cannot have all three of its named partners
three-regular.
\end{theorem}

Machine-checked. Were they, each would have this site as its unique live partner in its own
colour, so the complementary sums with the extra would vanish in \emph{every} colour --- not
merely one, which is all the degeneracy theorem gives. Support minimality still forces a
surviving colouring of that complement; its pinning puts a live row's named partner on its own
colour; and that is exactly what a three-regular named partner cannot survive, its only partner
in that colour being the excluded site.

The immediate consequence is that the excess sites cannot number two. With two, each has degree
exactly four --- at most three three-regular partners, and the other excess site as its extra
--- so all three of its named partners are three-regular, which the theorem forbids. With the
earlier result that excess degree never sits alone, the excess set is empty or has at least
three members.

What blocks the extension is precise: at three or more excess sites, a named partner may itself
be excess, and then only one colour is known degenerate rather than all three. Recovering the
argument needs degeneracy in the colours whose named partner carries excess.

\paragraph{The flattening has rank three.} One principle explains the cut theorems uniformly
and costs nothing to state.

\begin{theorem}[No rank-two flattening]\label{thm:rank2}
If for some bipartition the amplitude is a sum of two terms, each a function of the colouring
on one side times a function of the colouring on the other, the equations are contradictory.
\end{theorem}

Machine-checked, and with no tensor machinery: read at the two-sided constant colourings the
flattening is diagonal with non-zero entries, since a colouring constant on each side with
different constants is not monochromatic; a sum of two products has rank at most two, because
three vectors in a two-dimensional space are dependent; and pairing that dependency against the
diagonal kills a non-zero amplitude.

Connectivity is the case of no crossing edge, where the amplitude is a single product. The next
case is new.

\begin{theorem}[Two crossings at one site]\label{thm:twocross}
Let an odd part of the sites have all its outward live edges at a single site, reaching at most
two partners, with rank-one blocks. Then the equations are contradictory.
\end{theorem}

Machine-checked. Expanding at that site, the terms that keep it inside the part die on parity
--- the part minus two sites is still odd --- and two products remain. This strictly generalises
the bridge theorem, which is the one-partner case, and it is the first argument here that
consumes the \emph{rank} of a crossing block rather than only its support.

The general shape is worth recording: an odd part whose outward edges all sit at one site, with
rank-one blocks, needs at least three outward partners, since the colouring on the inside enters
only through that site's colour and the matching sum on the rest. Interior sites of such a part
have all their partners inside, so with minimum degree three the part has at least five sites
and its outward site is a cut vertex. The excess subgraph has no such cut, which is why this
lever, like the others, constrains the ambient structure without closing the remaining case.

\paragraph{A rank statement at every site.} The flip relation, contracted against a covector,
gives something the programme had not had: a rank constraint holding everywhere, with no
hypothesis on the support.

\begin{theorem}[Colour columns span]\label{thm:colspan}
A covector annihilating every colour-$k$ column at a site has zero $k$-component. Equivalently
$e_k$ lies in the span of the colour-$k$ columns, at every site of every solution.
\end{theorem}

Machine-checked. Pair the covector with the vector of complementary sums: each row's pairing is
the amplitude on the diagonal and zero off it, so the contraction is the covector's own
component times the amplitude --- and it is also zero term by term. This is the exact dual of
the column lemma: the fan says each colour has a column-supported partner; this says that
colour's columns reach that colour's direction.

Its first consequence is diagonality from support alone.

\begin{corollary}[Dead extra column]\label{cor:deadcol}
At a site whose live partners are its three named ones and a single extra, a dead colour-$k$
column of the extra block makes that colour's named block the single diagonal entry.
\end{corollary}

Every earlier route to diagonality went through a vanishing matching sum --- through deadness,
or through a degenerate colour. This one reads it off where the support is, and it costs
nothing.

That restructures what is left. The extra is live, so at least one of its columns survives, and
the remaining case splits by how many do: with one, the other two named blocks are diagonal and
the extra block is column-supported in that colour, so the site has two column-supported
partners there; with two or three, diagonality is free in the remaining colours and the
surviving columns are definite objects. The previous obstruction --- that only one colour could
be shown degenerate --- is no longer the operative one, since diagonality now comes from support
in every colour the extra does not reach.

\paragraph{Inert edges can be deleted.} The certificates all come from non-vanishing matching
sums, so the case where none exists has to be disposed of separately --- and it disposes of
itself. Suppose the matching sum on the complement of $\{u,v\}$ vanishes for \emph{every}
colouring. Then setting $W_{uv} = 0$ changes no amplitude: the term of the expansion at $u$
that uses the edge is multiplied by that sum, and no other term sees the edge, because every
other term's matching sum runs over a set not containing $u$. The modified system satisfies
the same equations with strictly fewer live pairs.

So one may assume, of every live edge, that some colouring of its complement has
non-vanishing matching sum --- and then the pinning theorem has something to bite on. This
is the lever for the case where all three colours are degenerate: there the constant
colourings give nothing, but a non-constant one must exist, and it is non-constant precisely
because the constant sums vanish, so span domination applies to it with no side condition.
This is machine-checked.

\paragraph{The fan relation.} All of the site-local statements above are readings of one
identity, and it is worth having in that form.

\begin{theorem}[The fan relation]\label{thm:fanrel}
At a site whose live partners are its three named ones and a single extra $v$, and for every
colour $\alpha$ and every choice of test vectors,
\[
  A_\alpha \prod_{w \neq u} \psi_w(\alpha)
    \;=\; \langle \mathrm{row}_\alpha W_{uv},\, \psi_v\rangle\, \hat H(\psi)
      \;+\; \sum_j W_{(u,\alpha)(w_j,j)}\, \psi_{w_j}(j)\, K_j(\psi),
\]
where $\hat H$ is the matching sum on the complement of $\{u,v\}$ and $K_j$ that on the
complement of $\{u,w_j\}$.
\end{theorem}

Machine-checked, with no hypothesis on the named blocks beyond column support. Setting every
$\psi_{w_j}(j)$ to zero kills the right-hand sum and the left side at once, and recovers the
support theorem; testing against basis vectors recovers the flip relation; testing at the
constant colourings recovers the constant reading.

It also recovers the pinning, in three lines instead of through separation and span
domination. Put a test vector missing colour $\alpha$ at the $\alpha$-named partner. The left
side dies, that coordinate being one of its factors; the named sum dies, diagonality leaving
only the $j = \alpha$ term and that carrying the same coordinate; and what remains is the
extra's row against a test vector times the matching sum on the complement. If the row is
live, the matching sum vanishes. Both the test-vector and the colouring form of this are
machine-checked, and the derivation needs no separation argument at all --- which also
retires the circularity that the old route carried.

\paragraph{A dead row leaves a monomial.} The fan relation reads the other way too. Put a
test vector concentrated on colour $\alpha$ at the $\alpha$-named partner. A dead row kills
the extra's term outright, diagonality leaves only the $\alpha$-named term, and what remains
determines the matching sum on the complement of $\{u,w_\alpha\}$ completely.

\begin{theorem}[Dead row, determined complement]\label{thm:deadrow}
If the named blocks are single diagonal entries and row $\alpha$ of the extra block is dead,
then the matching sum on the complement of $\{u, w_\alpha\}$ equals $A_\alpha/\lambda_\alpha$
at the constant colouring $\alpha$ and vanishes at every other colouring.
\end{theorem}

Machine-checked. A whole complement's matching sums collapse onto a single colour --- in
particular the sum vanishes for every colouring that sends the extra $v$ anywhere but
$\alpha$, and the colour-$\beta$ constant sums on that complement vanish for both
$\beta \neq \alpha$. This is the sharpest structural consequence yet extracted from a single
extra partner, and it is what a contradiction in the all-degenerate case has to be built
from: the surviving configurations are now pinned to a shape in which two of the three
complements are single monomials.

The determination has a consequence at the other end of the edge. Read the flip relation at
$w_\alpha$ instead of at $u$. The edge back to $u$ is monochromatic, so it contributes only
against row $\alpha$ --- and there it already supplies the whole colour-$\alpha$ amplitude,
the dead row having removed the extra from $u$'s own relation.

\begin{corollary}[The site carries its named partner]\label{cor:carried}
With a dead row $\alpha$, the aggregate of $w_\alpha$'s remaining colour-$\alpha$ edges is
zero against every row: $u$ alone accounts for $w_\alpha$'s colour-$\alpha$ amplitude.
\end{corollary}

Machine-checked. This is a constraint on a \emph{different} site's neighbourhood, which is
what the site-local arguments have been unable to reach.

\paragraph{The conjecture, modulo one statement.} Everything established assembles, in the
official formulation, into a single implication.

\begin{theorem}[Main, conditional]\label{thm:main}
If every GHZ system on more than four sites is three-regular, then no GHZ system in three
colours exists on more than four sites --- for every vertex count.
\end{theorem}

Machine-checked. Two points deserve emphasis. First, there is no base case: the six-vertex
statement is not used anywhere. It was carried as a sealed hypothesis while the uniform
argument was being built, and the uniform argument turns out not to need it, because the
spread case is settled for every even count above four rather than by descent to six.
Second, there is no finite check of any kind in the dependency cone, and no axiom beyond
propositional extensionality, choice and quotient soundness.

So the whole conjecture now rests on a single named hypothesis --- and that hypothesis says
less than it first appears. The column fan already supplies three live partners at every
site, one supported in each colour, and they are distinct. Three-regularity therefore makes
no claim about colour structure at all; it claims only that there is no \emph{fourth} live
partner. Loop weights, which appear in no matching and so in no amplitude, normalise away.

\begin{theorem}[Main, in support terms]\label{thm:maindeg}
If no site of a GHZ system on more than four sites has a fourth live partner, then no GHZ
system in three colours exists on more than four sites --- for every vertex count.
\end{theorem}

Machine-checked. What is open is a statement about the support of $W$ and nothing else.

\paragraph{The live graph is connected.} One global fact comes free and is worth recording.
Suppose the sites split into two non-empty parts with no live edge between them. Every
matching that crosses contributes zero, so each amplitude factors as the product of the two
parts' matching sums. Since every constant amplitude is non-zero, so is every part's constant
sum, in every colour. Now colour one part $k$ and the other $k' \neq k$: the colouring is not
constant, yet its amplitude is a product of two non-zero numbers. So no such split exists.

Both this and the factorisation it rests on are machine-checked. The factorisation is proved
by induction on the part's size rather than by a bijection between matchings: expand at one of
the part's sites, discard the crossing terms, and apply the induction to what is left. Stated
for any part of any ambient set, it is the first genuinely global tool in the development.

The same reading works for a \emph{live} row, at the cost of a condition. Choose test vectors
that annihilate the extra's row in colour $\alpha$ and miss colour $\alpha$ somewhere; the fan
relation again leaves only the $\alpha$-named term.

\begin{theorem}[The named partner, spared]\label{thm:namedspared}
The matching sum on the complement of $\{u, w_\alpha\}$ vanishes on the whole subspace where
the extra's colour-$\alpha$ row is annihilated and some site misses colour $\alpha$.
\end{theorem}

Machine-checked; with a dead row the annihilation is automatic and this is
Theorem~\ref{thm:deadrow}. Reading it as a divisibility: restricted to test vectors
annihilating the row, the complement's form is a single monomial times a linear functional at
the extra. Two of the three complements at a site are thereby pinned to monomials, and the
third to a monomial on a hyperplane.

\paragraph{A route that does not work.} It is tempting to try to bound the number of live
rows of the extra block by playing two of them against each other: if rows $\alpha$ and
$\beta$ are both live, $\hat H$ is divisible by $\psi_{w_\alpha}(\alpha)$ and by
$\psi_{w_\beta}(\beta)$, and one hopes to kill the left side of the fan relation while the
$j = \alpha$ term survives. It does not close. The left side's factors are the
$\psi_w(\alpha)$, all in colour $\alpha$; killing it through $w_\beta$ needs
$\psi_{w_\beta}(\alpha) = 0$, which is a different condition from
$\psi_{w_\beta}(\beta) = 0$, and imposing both only reaches colourings that send $w_\beta$ to
the third colour. Those never include the constant colouring the contradiction would need.
Killing the left side through $v$ instead is available, and yields that the matching sum on
the complement of $\{u,w_\alpha\}$ vanishes whenever the colouring sends $v$ off $\alpha$ and
$w_\beta$ off $\beta$ --- true but not contradictory, since the sum known to be non-zero
there is the one at the constant colouring $\alpha$, which sends $v$ to $\alpha$.

\begin{corollary}[Rank support]\label{cor:triple}
In the surviving case, if some test tuple on the complement of $\{u,v\}$ has non-vanishing
complementary sum, then $\operatorname{im} W_{uv}$ lies in the span of the other extra
partners' selected columns. In particular at least $\operatorname{rank} W_{uv}$ of those
columns are linearly independent, so $u$ has at least that many extra partners besides
$v$.
\end{corollary}

The three-fold conclusion --- three further extra partners with independent selected
columns --- is the full-rank subcase only. A block of rank one or two forces
correspondingly less, and the earlier statement, which omitted the rank hypothesis, was
wrong on that point.

No determinantal normal form is needed, and none is claimed. The containment above is
obtained constructively: given $y \in \operatorname{im} W_{uv}$ outside the span, separate
it by a covector annihilating the span --- ordinary finite-dimensional linear algebra ---
and the relation empties. That is safer than asserting membership in a determinantal
ideal, which would need the radicality of a pull-back that is not established here.

The three named partners cost nothing in this reading. For each colour $k$, either
$\rho_k = 0$, or the annihilation at $w_k$ forces $\psi_{w_k}(k) = 0$, since the named
block is $c_k \otimes e_k^{*}$ and so $z_{w_k} = \psi_{w_k}(k)\,c_k$; either way the $k$-th
term on the right vanishes. So no condition on the test vectors at the named partners has
to be imposed --- it comes free.

\subsection*{When the extra edges form a matching}

Suppose the extra edges form a perfect matching $E$, so every site has exactly one. Then
the fourth consequence above says no extra edge meets a constant colouring, and therefore
the $\gamma$-live graph has $M_\gamma$ as its \emph{only} perfect matching, for each
$\gamma$.

Now take a mixed matching $M$ inside the three named ones and colour by it. The matchings
compatible with that colouring are $M$ itself together with those replacing part of $M$ by
extra edges; and the part replaced must be a union of $M$-edges \emph{and} of $E$-edges,
hence a union of components of $M \cup E$ --- alternating cycles. Because the cycles are
disjoint the sum factorises, and the vanishing amplitude reads
\[
  \prod_{C \,\subseteq\, M \cup E} \bigl(1 + r_C\bigr) \;=\; 0,
  \qquad
  r_C \;=\; \frac{\prod_{e \in E \cap C} W_e}{\prod_{e \in M \cap C} w_e},
\]
the numerator read at the colours the colouring assigns. So some cycle has $r_C = -1$
exactly.

A cycle whose $M$-edges all carry one colour cannot: it is then a cycle of
$M_\gamma \cup E$, its colouring is the constant $\gamma$, and the uniqueness just
established forces one of its extra edges to be dead there, so $r_C = 0$. Hence:

\begin{proposition}\label{prop:noconst}
No mixed matching has its colouring constant on the extra edges.
\end{proposition}

That is a purely combinatorial demand, and contracting each extra edge makes it plainer.
A colouring constant on extra edges descends to the $n$ contracted nodes; the condition
that its parts be closed under the named matchings makes each named matching a
two-regular graph there, whose cycles partition the nodes. So the proposition says: the
three cycle systems on $n$ nodes admit no non-trivial partition into three parts, the
$\gamma$-th a union of $\gamma$-cycles. In particular no two of them have a disconnected
union.

That is the same shape as Lemma~\ref{lem:fourth} one level down --- but on half as many
nodes, and, crucially, it is \emph{false} there. On three nodes a two-regular loopless
graph must be a single triangle, so all three cycle systems are the trivial partition and
no non-trivial splitting exists at all. Three nodes is six sites, exactly the case sealed
off as a hypothesis; the failure is not an accident but the shadow of it. On four nodes it
can fail too: three four-cycles with connected pairwise unions admit no splitting either.

So the extra-matching case does not close on combinatorics alone. What
Proposition~\ref{prop:noconst} then says is a condition on the \emph{weights}: for every
mixed matching some alternating cycle has ratio exactly $-1$. That is one codimension-one
condition per mixed matching, and there are at least as many mixed matchings as
Lemma~\ref{lem:fourth} supplies. Whether that system is contradictory is the question this
case now reduces to --- and it is a question about weights again, not about graphs.

The reduction is still worth having: it replaces "some weighted matching sum fails to
vanish" with "a specific product of ratios equals $-1$, simultaneously across all mixed
matchings", which is far more rigid. But the earlier reading, that the case closes
combinatorially, was wrong, and is corrected here.\section{Theorem ledger}\label{sec:ledger}

\paragraph{Machine-checked, standard axioms only.} The pinning at a lone live neighbour,
off the diagonal and on it; that a well-behaved subset carries a solution, and the
restriction and relabelling that transport it; the site relation at any degree, the
vanishing it forces under independence, and the impossibility of independence everywhere;
that liveness forces degeneracy; that no plane holds every column at a site; the
resistance of an unused colour to
perturbation; the absence of an annihilator at a
site and the consequent failure of flat outer-product caps; the self-defeat of the
outer-product route where partners are few; the exclusion of rank-two blocks at low-degree
sites; amplitude naturality and matching
extraction; orbit normalisation and the classification of matching triples; the
four-vertex solution and the exclusion of colour-per-matching solutions above it; the
even-cycle two-colour solutions; the support restriction and star collapse for the
contraction; the decoupling of the two cap functionals and the explicit witness among
three vectors; the cap-incidence map for arbitrary weights and the rank
characterisation; the exact quadratic expansions at product and matching-sum level, the
second-order absorption identity, and the factorisation of its coefficient; the block
decomposition, the degeneration lemma, the dichotomy and its unconditional form under a
degree bound; live partners; the uniqueness of the live colour on colourings; and, at a fully degenerate site, the linear system solved by the complementary sums, the existence of a partner per colour, purity, the subsingleton support of a live edge, the fan of three distinct partners, the descent it yields under full support, its failure at eight sites, the symmetry of the complementary sums, the single colour carried by every participating edge when every site is degenerate; and, with no hypothesis beyond the equations, the collapse of the test sum, its vanishing under annihilation at every partner, the existence of a pure partner at every site in every colour; the splitting and factoring of sums over colourings at one site; the isolation of a single pairing and the compatibility it forces at both ends of an edge; the master relation and the column lemma it yields; the fan for a generic covector; that a colouring with exactly one contributing matching is constant; that a matching mixing two of three disjoint colour classes is refuted when the live edges are exactly those classes; that three live partners force each block to a single diagonal entry, that a three-regular configuration is three disjoint perfect matchings, and that such a configuration admits no mixed matching; that orthogonality localises the witness, of which the column lemma is the empty-set case; the support theorem for colourings of a complement; the general two-covector spared relation and the span-domination theorem it yields, together with the vanishing of a named partner column the colouring form in which the span is spanned by literal columns, and the pinning of the spared block live rows; the single-flip formula writing a one-site-flipped matching sum in single-colour ones, the flip orthogonality it yields at every site and colour, its confinement to two terms at a site with one extra partner, and the diagonality of a named block whose complementary sum vanishes; the vestigial colour pins in the paired sum, the spared relation for an arbitrary covector, span domination with no hypothesis on the named blocks, its reading at a constant colouring, the vanishing of the complementary sum in some colour at a site with one live extra; that an edge whose complement has vanishing matching sum in every colouring can be zeroed without changing any amplitude; the fan relation, the single identity of which every site-local statement above is a reading; the pinning read straight off it, in both the test-vector and the colouring form; the determination of a named partner's complement to a single monomial when the corresponding row of the extra block is dead; that the site then carries that partner's whole amplitude in the colour; the factorisation of a matching sum across a cut and the connectivity of the live graph that follows; that a perfect matching inside three disjoint ones is the same thing as a consistent colouring saying which each site uses, and that such a colouring exists non-constantly when two of the matchings generate a disconnected graph --- so that a three-regular solution with a disconnected pair is impossible; in the cycle coordinates, the fourth matching cut out by a chord across the two classes and the one cut out by two interleaved chords, the compatibility and non-constancy of each, the existence of an interleaved pair whenever every chord stays inside a class; the cycle coordinates themselves, built from any two disjoint involutions with connected union, and the disjointness of the two classes; hence a fourth matching for any three pairwise disjoint involutions of more than four sites, and hence that no three-regular solution exists above four sites; and the passage from the generic fan to column-supported blocks, together with the distinctness of the three partners it names.

\paragraph{Proved on paper, not yet formalised.} That a family of pairs fits inside one
star or triangle exactly when it is pairwise intersecting; and the connected case of the
fourth-matching lemma, which is now machine-checked in full. The consequence is recorded as
a single theorem: a three-regular solution of the equations on more than four sites is
impossible, for every site count and not merely for six.

\paragraph{Computation, with controls.} The pinned family across every cubic class at eight
sites, consistent and mismatched ends, against a four-site control that converges readily:
no start converges at eight. The dense frontier as well: complete bipartite, complete,
and complete-minus-a-matching supports at eight sites, with an analytic Jacobian whose
four-site control converges at every start --- no start converges at eight. These are
searches, and their silence is evidence about the formulations tried, not a theorem.

\paragraph{Computation, inconclusive.} The rank-one family at eight vertices. The search
is validated on both sides --- four vertices recovered, six vertices correctly refused ---
but at eight it neither finds a solution nor descends convincingly, and the gauge freedoms
were not quotiented.

\paragraph{Computation, with controls.} The least number of perfect matchings of a cubic
graph carrying three disjoint perfect matchings, enumerated at six, eight and ten
vertices: four, five and six respectively, never three. This tested the fourth-matching
lemma before it was proved, and the proof does not rest on it.

\paragraph{Computation, with controls.} Every four-site solution returned by the
validated search has single-entry edge blocks throughout, and the pure-partner conclusion
holds on each; random weights satisfy it in no trial. Evidence about four sites and about
the formulation searched, not a theorem about all sites.

\paragraph{Computation, exact arithmetic.} The infeasibility of the translation-invariant
family at eight vertices, with the search validated on the four-vertex solution; a cap on
the cube at eight vertices, all
ninety flatness constraints verified over the rationals; the absence of caps at
wide-range weights on the cube and on Wagner's graph; the coincidence of cap appearances
with singular edge blocks; the block decomposition checked at four and six vertices.

\paragraph{Conjectural.} That the weighted matching sum of a complement cannot vanish
identically on the product of the kernels --- equivalently that every live partner is
column-supported and its colour served once, which is what delivers the rigid shape; that the
solution variety meets the locus where an edge block is singular, at supports not covered by the degree bound; and that the colouring
condition in which the degeneration chain terminates cannot be met.

\paragraph{Not addressed.} Weights whose used columns span fully at every colouring. The
degeneration chain is vacuous there, and nothing here bears on them.

\paragraph{Withdrawn.} See Section~\ref{sec:withdrawn}.

\section{Forced contributions}\label{sec:forced}

A matching sum runs over very many matchings, so exhibiting one non-zero term proves
nothing: the others may cancel it. Every step that wants to \emph{produce} a
non-vanishing quantity has to defeat that cancellation. Two general tools now do.

\subsection*{The non-cancelling partner}

Let $\mathrm{pm}(c,S)\neq 0$ and let $u\in S$ be any site. Expanding at $u$ writes
the sum as $\sum_{v} W\!\left(u^{c(u)}v^{c(v)}\right)\,\mathrm{pm}(c,S\setminus\{u,v\})$.
If every summand had a vanishing factor the sum would vanish. Hence some partner $v$
carries \emph{both} a non-zero weight and a non-zero matching sum on what is left.

The second conjunct is the point. A live partner is easy to find and says nothing,
because the matchings through it may cancel; this returns one through which they
demonstrably do not. The construction repeats on the smaller set with a freshly chosen
site, so iterating builds a perfect matching all of whose tails are non-zero.

\subsection*{The product formula}

If no live edge joins two colour classes of a colouring $c$, then
\[
  A(c) \;=\; \prod_{k}\ \mathrm{pm}\!\left(c,\ c^{-1}(k)\right),
\]
by two applications of the cut factorisation. Support minimality supplies the
hypothesis: on a support-minimal representative every live pair is non-inert, hence
monochromatic once every site is degenerate.

This is the reformulation the equations want. The weights become three weighted graphs
$G_0,G_1,G_2$, one per colour --- no live edge mixes colours at its two ends, though a
pair may still carry weight in more than one colour --- and the system reads
\[
  \operatorname{haf}(A_k)\neq 0 \ \ \text{for each } k,
  \qquad
  \prod_k \operatorname{haf}\!\left(A_k[S_k]\right)=0 \ \ \text{for every non-constant partition.}
\]
No mixed amplitude survives to cancel against anything; the whole content sits in
hafnians of principal submatrices.

\subsection*{The equations that follow}

\emph{Two colours.} A set whose complement still matches in one colour carries no
matching sum in any other. Taking the set to be a single edge: a live edge annihilates
the other colours' complementary sums. Taking it to be a union of pairs gives the
four-site identities.

\emph{Colour exclusivity.} If a pair is live in a colour \emph{and} the rest still
matches in that same colour, the pair is dead in every other colour --- an edge that
mattered in two colours would annihilate its own complement. The colour classes are
therefore disjoint by proof rather than by assumption. Combined with the partner lemma:
every site has, in every colour, a non-cancelling partner exclusive to that colour, so
the minimum degree is three with no case analysis.

\emph{Disjoint pairs.} Two disjoint pairs of different colours cannot both be live over
a complement that still matches in the third.

\subsection*{Forced matchings, and the master contradiction}

The partner lemma manufactures a non-vanishing complement but leaves the choice of
partner to the construction. When a site has exactly \emph{one} live partner inside the
set under consideration, no choice is made at all: the expansion there has a single
term, and the two matching sums differ by a non-zero factor. This is the forced peel,
and it works on any set at any stage.

Two dual certificates follow. A \emph{forced peel} of $S$ takes it down to nothing, one
pair at a time, each seed having exactly one live partner left. A \emph{forced tail} at
$T$ peels the complement of $T$ away from the whole site set the same way. Neither
leaves room for cancellation.

They collide. A proper non-empty $S$ that peels forcedly to nothing in one colour and
forcedly away from everything in another cannot exist: the two-colour equation says a
set whose complement still matches in one colour carries no matching sum in another, and
both certificates say otherwise.

Both halves speak only about $S$, so the statement is usable at whatever scale the
configuration offers, with no global matching needed on either side. Two consequences
worth naming:

\emph{An alternating four-cycle is fatal.} Four sites carrying a cycle that alternates
between two colours, with the first colour's edges forcedly peelable and the second
colour absent from one diagonal, contradict the equations outright.

The corollary usually drawn from this --- that a colour class which is a perfect matching
has all its alternating cycles of length at least six --- is \emph{not} new: a shorter
cycle's vertex set is closed in both colours, which the two-colour result already forbids.
What the statement above adds is the case where the classes carry extra edges, so that the
vertex set is no longer closed and the uniqueness and diagonal hypotheses do the work
instead.

\emph{More generally}, an alternating cycle of any even length with no chord in the
second colour is an instance, whatever its length, provided it does not exhaust the
site set. That is the classical ``at least two cycles'' step, now uniform.

\emph{A unique partner cannot cancel.} If a site has one live partner in a colour, that
colour's complementary sum is automatically non-zero and the edge is exclusive to its
colour. So where every site has exactly one partner per colour the configuration is
three-regular, which is already excluded; and the entire residual case is a site with
two partners of the same colour. That is the only place a complementary sum is free to
vanish, and therefore the only place cancellation can hide.

\subsection*{The exact peel}

The forced peel asks every other partner of the peeled site to be dead. That is more than
the argument needs: it is enough that every other partner's \emph{term} vanish. A branch
dies when its weight is zero, and equally when what it leaves behind carries no matching
sum at all --- for instance when removing that partner strands some third site with
nothing to pair with.

The distinction is the difference between a site of degree one and a site whose other
branches lead nowhere, and the second is far more common. Two triangles joined by an edge
are the smallest example: no site has a unique partner, yet the perfect matching is
unique and every alternative strands a site.

Every forced peel is an exact peel, and the master contradiction runs on the stronger
certificate. The gain is exactly where the earlier version stopped: an alternating cycle
carrying a chord in the peeled colour is still exactly peelable, because the chord branch
strands the neighbour it leaves behind. A single fourth neighbour no longer blocks the
argument; what blocks it is a genuinely competing matching, which needs several.

\subsection*{Colour cuts}

The cheapest certificate in the argument needs no peeling at all. If no live edge of a
colour joins a set to its complement, that colour's matching sum factors as the product of
the two parts' sums. The whole one is non-zero, so both parts are --- with no choice made
and no cancellation to rule out. It applies to any union of connected components of a
colour class.

Two consequences follow immediately, each using only the equations: no minimality, no
degree bound, no case analysis.

\emph{The union of any two colour classes is connected.} A proper set left by neither
colour would carry the first colour's matching sum, and its complement the second's, which
the two-colour equation forbids. When the classes are perfect matchings this is the
classical statement that two of them union to a single alternating cycle through every
site. This is \emph{not} new here: the development already contained it, in a more general
form that assumes only a single mixed crossing condition rather than that every live edge be
monochromatic. The version above is a specialisation of that one, kept because the
monochromatic form is what the surrounding material discharges directly.

\emph{No non-constant colouring keeps every class inside its own colour.} Each class would
carry a non-zero matching sum, the product formula makes the amplitude their product, and a
non-constant colouring has amplitude zero. So every non-constant colouring has some class
from which its own colour escapes.

The second statement closes the three-regular configuration at once --- a second route to a
result the development already establishes by cycle coordinates, not a replacement for it:
there a colour class \emph{is} its matching, its components are its single edges, and the
mixed-colouring
construction supplies exactly a non-constant partition into unions of them.

What remains is therefore purely structural, with no hafnian and no cancellation in it:
whether there is a non-constant partition $V = S_0 \sqcup S_1 \sqcup S_2$ in which each
$S_m$ is a union of connected components of the colour class $G_m$. At degree three the
components are single edges and such partitions abound; at higher degree the components are
larger and the partitions harder to find. If every colour class is connected there are
none, and connected classes force many edges: each connected spanning subgraph needs at
least $n-1$ edges of its own colour.

\subsection*{The positive form}

The reduction the codebase carried all along asks that no support-minimal solution on more
than four sites have a site of degree four. That statement is vacuous once the conjecture
holds --- a GHZ configuration has an even number of sites, so any site set of more than
four carrying one is a case the conjecture covers --- and a vacuous hypothesis has no
non-trivial models. It can be refuted but never developed.

The same content states positively, and then it has objects in it:

\begin{quote}
On more than four sites, any weight system whose three constant amplitudes are non-zero has
some non-constant colouring with non-zero amplitude.
\end{quote}

The conjecture follows in three lines. The class quantified over --- weight systems with
non-vanishing constant amplitudes --- is full of concrete objects, and every certificate in
this development already has the shape the statement wants: each \emph{builds} a
non-vanishing matching sum rather than assuming one away. A set that peels exactly in one
colour whose complement peels exactly in another \emph{is} a mixed colouring with non-zero
amplitude; the two-colour equation was only how that fact got phrased as a contradiction.

\subsection*{The partner-pair dichotomy}

Working in that stance sharpens the four-site analysis. Take a live edge in one colour and a
partner for each of its ends in a second, with the complement of the four sites still
matching in that second colour. Then exactly one of two things holds:

\begin{itemize}
\item the pair of partners is dead in \emph{both} other colours, live only in the colour that
produced it if at all; or
\item the four sites carry a full four-cycle in the edge's own colour, so both ends of the
edge have a second partner there.
\end{itemize}

The third colour dies unconditionally. The edge's own colour dies unless the four-cycle
closes --- and a closed four-cycle is exactly the configuration in which the four-site
matching sum cancels rather than vanishing termwise. The dichotomy is therefore ``termwise
vanishing, or cancellation'', localised to four sites and made explicit.

Two negative results bound what can be hoped for here. The hafnian has no adjugate: the
off-diagonal analogue of the flip relation is false, so there is no matrix identity relating
a colour class to its complementary sums. And hafnians satisfy no Pl\"ucker relation --- an
exact check on six sites shows both sign conventions fail, the unsigned one by exactly the
terms that split four sites two-and-two across the complement, which is the double-counting a
Pfaffian's signs cancel and a hafnian's do not. Cancellation has to be defeated by
construction, which is what the peels, cuts, and parity killers do.

\subsection*{The crux, sharpened}

The last equation is the one that would finish the argument, and exactly one quantity
stands in the way. Let $e$ be a live $j$-edge and $f$ a disjoint live $j'$-edge with
$j\neq j'$. If
\[
  \mathrm{pm}\!\left(\text{const } k,\ V\setminus(e\cup f)\right)\ \neq\ 0
  \qquad (k\notin\{j,j'\}),
\]
the two live weights multiply to zero and the system collapses. Both weights are
non-zero by construction --- take the exclusive partners of two distinct sites --- so
the entire remaining content of the conjecture is the non-vanishing of that one
complementary sum.

The two-site version of it \emph{does} vanish: that is the live-edge equation. The
partner lemma manufactures non-vanishing complements, but only of unions of $k$-edges,
because it peels one colour at a time; $e$ and $f$ wear the other two. This mismatch is
the whole gap, and it is now a single statement about one matching sum rather than a
case analysis over configurations.

\subsection*{Cuts}

A separate line, and the only one here that does not need the live edges to be
monochromatic: it constrains how the live graph may be cut, and applies to every
configuration.

If a part of \emph{even} size is joined to the rest by a single live pair, the matching sum
still factors. Expanding at that pair's inside end, the crossing branch strands the part with
an odd number of sites and no way out, so it contributes nothing, and every other branch
removes the pair's inside end and leaves no crossing at all. Parity does the work through the
odd-part killer; no crossing-count machinery is needed. The amplitude then factors as though
the cut were clean, which gives it flattening rank one where the equations demand three. So
\emph{an even part joined by a single live pair is impossible}, with no minimality, no degree
bound, and no degeneracy assumption.

The odd companion is not the same statement. With the part odd and one live crossing pair,
every matching must use that pair, so the amplitude is one weight times the two sides' matching
sums --- but the crossing weight couples the two sides, so the flattening rank is bounded only
by the block's own rank, which may be three. The contradiction needs a rank-one block. In the
monochromatic setting that is automatic, every block being single-entry.

Two crossings behave the same way once the splits are available on an arbitrary site set.
Expanding at the first crossing's inside end sends that branch to the odd one-crossing split
and every other branch to the even one --- except the branch that removes the \emph{second}
crossing's inside end, where no crossing survives and the plain split applies. Both terms then
factor across the cut given rank-one blocks, so the amplitude has rank two and the equations
are contradicted. Together:

\begin{quote}
Every cut whose crossing blocks have rank one is crossed at least three times.
\end{quote}

The method stops there, and not by accident. With three crossing pairs the flattening bound is
four while the equations forbid only two; and the GHZ tensor's flattening rank across any
bipartition is exactly three, so a cut carrying three independent channels is what the
equations demand rather than what they forbid.

What this does not give is a bound on degrees. A single site is a cut, so the conclusion
recovers minimum degree three and says nothing about a site with a fourth neighbour.

\section{Pendant certificates}\label{sec:pendant}

The cuts of the previous section bound how a set may be joined to its complement. What
follows bounds nothing: it \emph{produces} non-vanishing amplitudes, and does so without
any cancellation argument at all.

\subsection*{Splitting along an invariant set}

A matching of a set is carried by a fixed-point-free involution. If a subset is closed
under that involution the involution restricts to it --- and, because the complement of
an invariant set is invariant too, it restricts to the complement in the same breath. So
one matching splits along any invariant set into matchings of \emph{both} sides.

Two matchings therefore split together: a set closed under both is matchable in the first
colour and its complement in the second. The two-colour picture needs no cyclic order and
no vertex count, only that a pair of involutions has a proper invariant set. What it does
not give is non-vanishing: the matching just built is one term of its class's sum.

\subsection*{The residual with no quantum content}

Proving the positive form by contradiction supplies the vanishing condition --- that is
what a contradiction hypothesis is --- and on a support-minimal representative with every
site degenerate it supplies monochromaticity too. The product formula then turns an
amplitude into a product of its classes' own-colour matching sums, and the converse holds:
if each class carries a non-vanishing matching sum in its own colour, the amplitude does
not vanish.

So the conjecture reduces on this branch to a statement about weighted graphs:

\begin{quote}
Three weightings on a common vertex set, each with a non-vanishing perfect matching sum,
always admit a non-trivial partition of the vertices into three parts, the $m$-th of which
has a non-vanishing matching sum in the $m$-th weighting.
\end{quote}

No amplitudes, no edge colourings, no quantum content --- and no vertex count.

\subsection*{Two pendant edges}

Using all three colours removes the cancellation argument entirely. Give one colour a
single live edge, a second colour another single live edge, and the third colour
everything else. A two-element class has exactly one matching, so its sum \emph{is} a
single weight; there is nothing to cancel. Only the third class is a genuine sum.

The resulting certificate needs no degree bound, no vertex count, and no estimate. Its
one substantial hypothesis is that the four sites carried off have a certified complement.

\subsection*{The colour symmetry}

Nothing in the equations distinguishes the three colours. Permuting the colours of the
weights preserves both the vanishing condition and the absence of live edges mixing
colours, and the amplitude of the permuted weights at a colouring is the amplitude of the
original at the permuted colouring. Every certificate proved for one assignment therefore
holds for all six, and is transported rather than restated.

\subsection*{Why four sites are exceptional}

At four sites the certified-complement hypothesis is free: the complement is empty and the
empty matching sum is one. The certificate then reduces to a bare statement --- no pair
live in one colour has a disjoint pair live in another --- and that is exactly what the
four-site solution achieves.

It achieves it for a reason that cannot survive. There the complement of a pair \emph{is}
a pair, and it wears the pair's own colour; for the exclusive-or colouring of four
vertices this is the identity $0\oplus 1\oplus 2\oplus 3 = 0$. Above four sites the
complement of a pair is not a pair, the mechanism has nothing to say, and the content
moves entirely into the certified complement.

\subsection*{The certified complement always exists}

And that hypothesis is always satisfiable. Expand the colour-$0$ sum at a site $a$ and
then at its colour-$1$ partner $p$. If $\{a,p\}$ is \emph{certified} in colour $1$ then it
is dead in colour $0$, so $p$ can never turn up as $a$'s own colour-$0$ partner and the
expansion at $a$ lands strictly elsewhere. Two peels then produce four distinct sites
$a,p,v,w$ with $W_0(a,v)\neq 0$, $W_0(p,w)\neq 0$, and a certified complement.

\subsection*{What remains, exactly}

The certificate always has its four-set. What it lacks is control over the two sites that
come off, and the equations force precisely those to be dead to each other in the third
colour:

\begin{quote}
Through any pair $\{a,p\}$ certified in colour $1$, the colour-$0$ expansion at $a$ and at
$p$ produces surviving pairs $(v,w)$; the equations say none of them is live in colour $2$.
\end{quote}

Every other hypothesis of every certificate above is free. The residual is a question of
steering an expansion onto a live pair --- fully explicit, with no cancellation estimate,
no degree bound, no vertex count, and nothing hidden in a choice left unmade.

\section{The certified colour graphs}\label{sec:certified}

A pair is \emph{certified} in a colour when it is live in that colour and its complement still
carries a matching sum in the same one. Certified pairs are what the peel produces, and they are
the right object for the rest of the argument.

\subsection*{Disjointness, and a degree of three}

A certified pair is dead in every other colour: were it live in another, the complementary sum
that survives would be annihilated. So the three certified colour graphs are pairwise
edge-disjoint. Every site has a certified partner in every colour, so their union has minimum
degree three.

The qualifier matters. This says nothing about a merely live pair, which may carry weight in two
colours; what is excluded is a pair that is live \emph{and} leaves its complement matching.

\subsection*{On a degenerate site every live pair is certified}

Support minimality makes a live pair non-inert, so some colouring survives on its complement.
Degeneracy at the site kills every colouring that is non-constant there. What survives is
therefore constant --- and a constant survivor is exactly a certificate.

So on the branch where every site is degenerate there are no uncertified live pairs at all: the
live graph \emph{is} the union of the three certified colour graphs, and since a certified pair is
dead in the other colours, a live pair wears exactly one. The colour classes are edge-disjoint,
which earlier in this programme could only be assumed.

\subsection*{The colour-degree equivalence}

The degree machinery counts live partners; the certified graphs are what the peel produces. On a
degenerate site the two neighbourhoods coincide, so each is a statement about the other. In
particular the live degree is the sum of the three colour degrees, each at least one --- a
decomposition that previously needed exclusivity as a hypothesis and now does not.

The bound between them, formerly available in one direction only, becomes an equivalence there:
live degree at most three if and only if exactly one live partner in each colour. So a second
partner in a single colour and a fourth live partner are the same condition, rather than one
implying the other.

\subsection*{What is left}

Everything on this branch assembles. If no site has a second live partner in any one colour, the
certified classes are perfect matchings, their union is cubic, and above four sites a cubic graph
with a one-factorisation admits a fourth perfect matching --- which is the contradiction the
three-regular exclusion supplies. Four vertices is exactly where that fourth matching fails to
exist, and it is the only place in the whole development where the vertex count does structural
work.

That is worth stating plainly as a constraint on any future step. Every certificate in
\S\ref{sec:pendant} is consistent with the four-site solution, because none of them uses the
vertex bound for anything beyond knowing that a four-element set is not everything. A card-blind
tool cannot force what the four-site solution satisfies, so no accumulation of them will close the
conjecture. What is needed is a fact true above four and false at four, of the fourth-matching
kind.

\subsection*{In the official formulation}

The library's solvability predicate is identified with the official equation system, so each
conditional result can be stated where the conjecture is actually posed rather than in a
paraphrase of it. Three such statements are recorded: from the colour-degree hypothesis, from the
partition question, and from the degree bound alone. None of them uses the six-vertex base case;
the uniform argument does not need it.

\section{The fibre criterion}\label{sec:fibre}

Everything in this section rests on one identity.  A colouring of the sites partitions them into
three fibres.  Where no live pair mixes colours, every matching contributing to that colouring's
amplitude must join sites of equal colour, so it decomposes into a matching of each fibre, and the
amplitude is the \emph{product} of the three fibres' monochromatic matching sums.  A non-constant
colouring has amplitude zero.  Hence:

\begin{quote}
If some non-constant colouring has all three fibre sums non-zero, the configuration cannot exist.
\end{quote}

Call a fibre \emph{rigid} when its own colour's matching sum over it does not vanish.  The criterion
asks for a non-constant colouring with three rigid fibres.

\subsection*{Where cancellation must live}

The first thing the criterion settles is the shape of the obstruction.  Cancellation cannot be
spread across colours: it must occur inside a single fibre, between two monochromatic matchings of
the same colour on the same sites.  Two such matchings differ by an alternating cycle, necessarily
monochromatic, and a monochromatic cycle requires every site on it to carry a second live partner in
that colour.  So a configuration surviving the criterion must carry excess, and the excess must be
concentrated enough to close a cycle within a single colour.

\subsection*{Rigidity from a matching}

The cheapest source of rigid fibres is a perfect matching.  Colour each site by the colour of its own
matched pair; the fibres are then the sites matched in each colour.  If, inside each fibre, no site
has a second live partner of that fibre's colour, the fibre's live graph in its own colour is the
matching itself, and the fibre is rigid.  A matching using more than one colour therefore kills the
configuration.

Purity is needed only inside a fibre.  A site carrying several partners in one colour is harmless
unless they are matched into the same fibre as itself, and the matching is ours to choose.

Where nothing is in excess, every pair is pure and the criterion reduces to the existence of a
perfect matching of the three colour classes using more than one colour.  That is elementary.  If two
colour classes meet in two or more cycles, recolouring one of them gives such a matching.  Otherwise
every pair of classes meets in a single Hamiltonian cycle, and the third class is a set of chords.  A
chord joining sites at odd distance leaves two arcs of even size, each matched by its own cycle
edges; with the chord that is a mixed perfect matching.  If instead every chord stays inside one
parity class of the cycle, then an even-class chord crossing an odd-class chord leaves four arcs of
even size, and two chords with those arcs again give a mixed perfect matching.  Crossing cannot be
avoided: take the chord whose arc is shortest, of length $d$; because the classes alternate around
the cycle, that arc contains $d/2 \ge 1$ points of the other class, which no chord may cross, so they
are matched among themselves by a chord of arc length at most $d-2$.

At four sites the argument correctly fails.  There the only perfect matchings are the three colour
classes themselves, each of them constant, and that is exactly how the four-site solution survives.

\subsection*{Rigidity for free}

Rigidity need not be checked at all when a fibre is closed in its own colour, meaning that no live
edge of that colour leaves it.  Then the colour's matching sum factors across the cut into two
factors whose product is the colour's amplitude, which does not vanish; so neither factor does.

A set closed in a colour is a union of connected components of that colour's live graph.  In this
form the criterion carries no arithmetic whatever: it asks for a partition of the sites into a union
of components of the first colour's graph, a union of components of the second's, and a union of
components of the third's, with not everything in one part.

Two colours never suffice.  If a colouring used only two, one part would be closed in its own colour
and its complement closed in the other, making a single set closed in two colours, and no
configuration admits one.  Consequently a colour whose live graph is connected can supply no fibre at
all, and the closed form of the criterion applies precisely when the three colour graphs are
disconnected and their components tile the sites.

\subsection*{One closed set is enough}

A tiling is more than is needed.  A single set closed in one colour already supplies all three
fibres: the set itself, split between the other two colours, and its complement, whose rigidity in
the closing colour comes free from the factorisation above.  Formally, if $A$ is closed in colour
$a$, both $A$ and its complement are non-empty, and $A = B \sqcup C$ with the $b$-sum over $B$ and the
$c$-sum over $C$ both non-zero, the configuration cannot exist.

The smallest instance is four sites: a set closed in one colour splitting into a live pair of each of
the other two.  The Hamiltonian construction above is another instance, with $A$ the complement of a
pure pair and $B, C$ the two arcs.

This form exposes a tension the earlier ones hide.  Extra live pairs make the split easier to find
and the closure harder to achieve.  A configuration blocking the criterion must therefore be
simultaneously rich enough to connect every colour graph -- which takes at least $3(n-1)$ extra live
pairs, an average live degree near six -- and poor enough to offer no closed set that splits.

\subsection*{The split, forced}

The criterion's remaining content can be isolated completely.  Take the fibre of the first colour to
be a set $F$ whose own-colour matching sum does not vanish -- a single live pair will do, its sum
being that pair's weight.  What must be produced is a division of the remaining sites into a part
rigid in the second colour and a part rigid in the third.

Rigidity of a part in a colour means that colour's live graph, restricted to the part, is exactly a
perfect matching of it.  So a site in the $b$-part has its $b$-partner in the $b$-part, and a site in
the $c$-part has its $c$-partner in the $c$-part.  Each rule is symmetric within its pair, so the
label is constant on every named $b$-pair and on every named $c$-pair, hence on each connected
component of the union of the two colour classes restricted to the sites outside $F$.

That union is a disjoint set of even cycles, and removing $F$ turns the cycles meeting it into paths.
The labels at the ends of those paths are forced: a site whose $b$-partner has been removed cannot be
in the $b$-part, so it lies in the $c$-part, and symmetrically.

Two consequences follow at once.  If $F$ is a single pair whose sites lie on \emph{different} cycles,
then removing one of them leaves its cycle as a single path whose two ends are its $b$-partner and its
$c$-partner -- one forced into each part, in one component.  That is impossible, so the two sites lie
on a common cycle.  And along that cycle the two arcs they leave must each begin and end with an edge
of the same colour, which happens exactly when the arc has an even number of sites, that is when the
two sites sit at odd distance.

Cycles meeting $F$ nowhere survive intact, and a whole cycle is a union of $b$-pairs and equally a
union of $c$-pairs, so each may take either label freely.

\subsection*{The case with nothing in excess, and what it leaves}

Here the rigidity conditions are automatic -- each colour's live graph is its named matching -- and only
the parity survives.  If two colour classes meet in several cycles, recolouring one already gives a
mixed matching and the criterion applies.  Otherwise they meet in a single cycle through all sites,
and the third class is a set of chords.

A chord joining sites at odd distance is exactly the pair required, and the construction closes: the
two edges at each of its endpoints carry different colours, so the two arcs receive different labels
and all three fibres are non-empty.

If every chord joins sites at even distance the construction is not available, and an attempt to use
two crossing chords instead fails.  Deleting four sites leaves arcs whose labels are again forced by
the colour of the edge joining each arc to a deleted site, and nothing makes those labels differ.  On
the eight-site cycle with its four diagonals, coloured alternately along the cycle with the diagonals
third, the two crossing diagonals leave two arcs which both receive the same label; the third fibre
is empty, the colouring uses two colours, and a set closed in one colour with complement closed in
another does not exist.  So that case is not settled by this construction.

That configuration is nevertheless covered, by the first case rather than the last: two of its colour
classes meet in two four-cycles rather than one, so recolouring one of them already gives a mixed
matching.

What the construction does not reach is the case where all three pair-unions are single cycles
\emph{and} every chord joins sites at even distance in each of them.  Such colourings exist.  Saying
that a colour's chords all join sites at even distance says that colour preserves the two parity
classes of the other two colours' cycle; taking the three parity functions together gives each site a
triple, on which the three colours act by adding the three even-weight vectors.  The skeleton is
connected, so the triple takes exactly four values and the three colours induce on those classes the
three perfect matchings of the four-site configuration: every colouring in this case is a covering of
it.

The covering is described by a permutation of the fibre on each edge, and a pair-union is a single
cycle exactly when the corresponding holonomy is a full cycle of the fibre.  Gauging the three edges
of a spanning tree to the identity leaves three free permutations, and the three holonomies work out
to $yz^{-1}$, $xz$ and $x^{-1}y$.  For a fibre of two elements a full cycle is the transposition, so
two of the holonomies agree and the third is forced to the identity -- impossible, which is why eight
sites admit no such colouring.  For a fibre of three the constraint is satisfiable, and the resulting
twelve-site colouring meets every condition of the case.

What this limits is the elementary \emph{existence} argument, not the criterion.  Where nothing is in
excess the criterion needs only a perfect matching using more than one colour, and such a matching
exists in every cubic graph carrying a proper three-edge-colouring except the one on four sites --
classically, because such graphs are bridgeless and only the four-site one has just three perfect
matchings.  The twelve-site covering above is no exception: it carries perfect matchings using all
three colours, so the criterion applies to it as it does to every other configuration without excess.

So the case with nothing in excess is settled, and settled by the criterion.  What the arc
construction supplies is an elementary proof of the matching's existence in two of the three
situations -- a pair of classes meeting in several cycles, or a chord joining sites at odd distance --
and the remaining situation is left to the classical fact.

At four sites every step still runs except the last: there the two classes form a cycle with a single
pair of chords, at even distance and not crossing, and indeed the only perfect matchings are the three
colour classes, each of them constant.  The exception is produced by the argument rather than assumed
by it.

\subsection*{What excess changes}

Excess adds live pairs beyond the named ones.  It enlarges the supply of first-colour fibres, since
any set with non-vanishing own-colour sum will serve, and it enlarges the supply of pairs at odd
distance.  What it costs is the rigidity clause: an arc labelled with a colour must contain no further
live pair of that colour.

The free cycles are where the room is.  Each may be labelled either way, so the obstruction is a
two-satisfiability question over their labels with the two arcs held fixed.  That is the open case,
and it is the only one.

\subsection*{What is not claimed}

None of these forms is a reduction.  A non-constant colouring's amplitude \emph{is} the product of its
fibre sums, so no genuine configuration admits three rigid fibres, and each criterion is equivalent
to the conjecture rather than weaker than it.  Their value is that they fix the shape of a
certificate and move the whole difficulty into a construction over the support: which pairs are
live, which colours they carry, and how the resulting three graphs are connected.

\section{The problem as one polynomial}\label{sec:poly}

Attach a variable $x_{v,a}$ to each site and colour, and form the matrix whose entry at a pair of
sites is $\sum_a W_a(u,v)\,x_{u,a}x_{v,a}$.  Its hafnian expands over matchings, and each matched pair
forces one colour at both of its ends, so

\[
  \operatorname{haf} M(x)\;=\;\sum_{c} \operatorname{amplitude}(c)\;\prod_{v} x_{v,c(v)} .
\]

The equations say exactly that this polynomial is supported on the three monochromatic monomials:

\begin{quote}
$\operatorname{haf} M(x) = \sum_{j} \lambda_j \prod_v x_{v,j}$ with every $\lambda_j$ non-zero.
\end{quote}

The conjecture is that this forces at most four sites.  Everything in this programme is a way of
reading that one identity.

\subsection*{What the earlier lanes are, in this language}

Refusing one colour per site and merging the other two is evaluation at a zero-one point: setting
$x_{v,a}=0$ exactly when $a$ is the refused colour leaves $\sum_{j \notin \operatorname{im} e}
\lambda_j$, which is why an onto refusal gives zero and a refusal missing one colour gives that
colour's amplitude.  The row identities obtained from a live pair are differentiation in a single
variable.  The rescaling that fixes the equations is the substitution $x_{v,a} \mapsto \mu(v,a)
x_{v,a}$, and the cube-root twists are the special case $\mu(v,a) = \omega^{a k_v}$.

That observation is worth stating because it settles the reach of all of them at once:

\begin{quote}
An identity obtained by evaluating or differentiating the polynomial identity is a consequence of the
equations, and cannot constrain the weights beyond them.
\end{quote}

So no amount of cleverness with test points, characters or derivatives closes the question by itself.
Each of those lanes was pushed to its limit and stopped, and this is the single reason.

\subsection*{What must be combined with it}

The coefficient of a colouring is the sum, over the matchings compatible with it, of their weight
products.  A colouring whose compatible matching is \emph{unique} therefore has a non-zero
coefficient, and must be constant.  Reading that contrapositively:

\begin{quote}
For every non-constant colouring, the live pairs admit either no compatible perfect matching, or at
least two, which cancel.
\end{quote}

A perfect matching of the live pairs induces a colouring, compatible by construction, and that
colouring is non-constant exactly when the matching uses more than one colour.  So a configuration
must either have no matching using more than one colour -- which is what happens at four sites, where
the only matchings are the colour classes themselves -- or every such matching must be cancelled by
another compatible with the same colouring.

Two matchings compatible with one colouring differ inside a single colour class, hence by a cycle
alternating between them in that one colour.  A monochromatic cycle needs every site on it to carry a
second partner in that colour.  So cancellation requires excess, concentrated in one colour along a
cycle.

This is where the difficulty now sits, and it is a statement about which pairs are live rather than
about their weights.

\section{Withdrawn claims}\label{sec:withdrawn}

\paragraph{Contraction fails on cubic supports from eight vertices upward.}
Withdrawn. The verdict came from a recognizer that tested each kernel basis vector for
a nonzero diagonal and a nonzero contraction scalar simultaneously. Those are two
linear functionals on the kernel and a witness may be a sum of basis vectors, so the
recognizer under-reported. An exact rational recognizer exhibits a cap on the cube,
verified in full. Everything built on the withdrawn premise falls with it: that the
route ``misses by two pairs'', that the obstruction is one unit of matching number,
that passing to second order is forced, and that the four-vertex matching sum is shown
to control every route.

\paragraph{Support-criterion failure implies no cap.} Withdrawn. The support
restriction is sufficient for an incidence entry to vanish, never necessary;
cancellation inside the cap matrix is not excluded by it. Censuses and degree tables
computed from the live-pair criterion describe that criterion's reach and not cap
existence.

\paragraph{Named cubic graphs settle the cubic case.} Withdrawn as stated. The cube,
Wagner's graph and Petersen are evidence about themselves; they were never a universal
theorem about cubic supports.

\paragraph{The two-colour cycle as evidence at three colours.} Withdrawn. It is a
cross-dimensional control on the recognizer only, and says nothing about the
hypothetical three-colour solution variety.


\begin{enumerate}
\item \emph{That descent is the conjecture.} The reduction ``every solution on $2n$
vertices yields one on $2n-2$, plus a base case, implies the conjecture'' is proved.
It was described as reducing the conjecture to one hypothesis. That description is
wrong: descent is strictly stronger than needed and may fail even if the conjecture
holds, since it demands descent of solutions no proof need ever meet.

\item \emph{That every vertex needs three partners.} Independence of the
monochromatic row is stated at one background colour, and there only that row is
independent of the other two --- the other two may be dependent, and at a vertex of
degree two must be. Changing the background repaints the columns, so rows at
different backgrounds are different vectors and do not combine.
\end{enumerate}

\section{The obstructions, precisely}\label{sec:obstruct}

\begin{openq}[A non-vacuous degeneracy]
The constant background forces the vertex's vector nonzero but yields only two
annihilations; the non-constant background yields three but no nonzero coordinate.
They can be combined --- recolour the vertex $v_0$ where the constant background is
nonzero; by locality that coordinate is a matching sum over a set excluding $v_0$
and does not notice.

\emph{The combination can be empty.} Recolouring $v_0$ moves the row entry at
$v_0$. Tested against the four-vertex solution above, all three rows vanish at that
coordinate in all twenty-four instances, so the annihilation holds for free. The
tension is structural: locality preserves the coordinate only because the recoloured
vertex is excluded, and that is the same vertex whose row entry is disturbed.

\emph{What is wanted:} a coordinate that is nonzero and at which some row survives.
\end{openq}

\begin{openq}[Unavoidability]
Local reductions are worth nothing without it. No argument here shows a minimal
counterexample must contain any configuration that has been discharged. This is
where the difficulty of discharging arguments has always lived.
\end{openq}

\section{Method}

Two habits produced most of what is above and both are worth stating.

\emph{Choose the representation that already knows the fact.} The cycle argument
stalled over $\mathrm{Fin}\,m$, where each statement about a successor needed a case
split on wrapping; over $\mathbb{Z}/m$ it went through as first written. Parity as a
ring homomorphism rather than a remainder made the alternation induction disappear
into \texttt{map\_add}.

\emph{Test a claim against a proved instance before building on it.} The
four-vertex solution serves as an oracle: any assertion about arbitrary solutions
can be checked against it first. This caught the vacuousness above before a layer
was built on emptiness, and caught a false hypothesis in a general theorem before it
was formalised.

\end{document}
